

\documentclass[a4paper,12pt]{article}
\usepackage[utf8]{inputenc}
\usepackage{multicol}
\usepackage{url}
\usepackage[numbers]{natbib}
\usepackage{multirow}
\usepackage{graphicx}
\usepackage{array}
\usepackage{makecell}
\usepackage{adjustbox}
\usepackage{seqsplit}
\usepackage{amsmath}
\usepackage{algorithm}
\usepackage{algpseudocode}
\usepackage{placeins}
% Redefinir @ para uso en direcciones de correo
\makeatletter
\newcommand{\email}[1]{\ttfamily#1}
\makeatother
% Language setting
% Replace `english' with e.g. `spanish' to change the document language
\usepackage[english]{babel}
\usepackage{multicol}
\usepackage{enumitem}
% Set page size and margins
% Replace `letterpaper' with `a4paper' for UK/EU standard size
\usepackage[letterpaper,top=2cm,bottom=2cm,left=2.5cm,right=2.5cm,marginparwidth=1.75cm]{geometry}

%Package for figures
\usepackage{tikz}
\usepackage{caption}
\usepackage [numbers]{natbib} % Use the numbers style
\usepackage{ragged2e}
\usepackage{float} % for the [H] specifier
\usepackage{marvosym}
\usepackage{wasysym}
\usepackage{amsmath}
\usepackage{graphicx}
\usepackage[font=scriptsize,labelfont=bf]{caption} % Adjusts the font size for captions
\usepackage[colorlinks=true, allcolors=blue]{hyperref}
\usepackage{makecell}
\usepackage{subcaption} % for subfigures
\usepackage{booktabs}       
\captionsetup[sub]{justification=centering}

\begin{document}
\date{}
\begin{titlepage}
	\begin{center}
		\vspace*{1cm}
		\includegraphics[width=\textwidth]{CRTD_TUD_logo.jpg}
		\vfill
		
		\huge{\textit{Team Project report}}
		
		\line(1,0){400}  \\ [2 mm]
		\huge{\textbf{
         Protein Ligand Interaction Guided Molecular Generation Using ChemBERTa}
		\line(1,0){400} \\ [5mm]
        \large{Summer term, September 2026}\break
        \small{\\Master of Science  program “Computational Modeling and Simulation”, Technical University of Dresden}\break
		\large{ \\ Mohammed Hasin Ishrak (5203084)\\}
        \large{ \\ Zaeem Asghar (5274291)\\}
        \large{ \\ Mahzabeen Chowdhury (5192937)\\}
		\vfill
		Advisors: Prof. Dr. Michael Schroeder, Philipp Schake, MSc
		\vfill
		
		
		\large{Dresden, Germany} \\
	\end{center}
\end{titlepage}
}
\clearpage

% Sections
%\tableofcontents
%\newpage
\begin{abstract}

Identifying novel small-molecule drug candidates with predictable target
engagement remains a core challenge in early-stage drug discovery, where
the vast combinatorial size of chemical space far exceeds the reach of
experimental screening. Transformer-based molecular language models such
as ChemBERTa learn rich atom-level representations from SMILES sequences,
yet their standard random-masking objective treats all atomic positions
as equally informative, rendering the reconstruction process blind to the
three-dimensional non-covalent interactions that govern binding affinity
at a protein active site.
Here we present GenPLIP, a zero-shot molecular generation pipeline that
couples per-atom interaction annotations produced by the Protein--Ligand
Interaction Profiler (PLIP) with ChemBERTa's masked-language-model
reconstruction mechanism through an interaction-aware
Non-Interacting-Part (NIP) masking strategy where atoms not participating in
pharmacophoric contacts such as hydrogen bonds, water bridges, hydrophobic
contacts, and $\pi$-stacking are selectively masked, directing the
model to regenerate peripheral scaffolds while preserving the
binding-critical core as unmasked context.
Benchmarked across five BRD4 co-crystal ligand systems, NIP masking
consistently generated greater libraries of structurally novel chemotypes
(Tanimoto $< 0.40$ to the reference ligand), anchored output to
pharmacologically validated ChEMBL analogues in the medium-similarity
range ($T = 0.40$--$0.70$), and produced top candidates with strong
predicted BRD4 binding (GNINA Vinardo $\leq -9.06\,\text{kcal mol}^{-1}$,
CNNaffinity $= 7.59$) that spontaneously converged on the para-chlorophenyl
and triazole pharmacophore essential for acetyl-lysine mimicry without
any explicit pharmacophore constraint or task-specific fine-tuning.
These findings demonstrate that structure-derived non-covalent interaction
priors can be encoded into a pre-trained molecular language model
through selective masking alone, providing a tractable zero-shot
pathwaythat bridges sequence-based deep learning and structure-based
drug design.
\end{abstract}
 
\noindent\rule{\linewidth}{0.5pt}
\vspace{0.5em}


\section{Introduction}
The process of discovering new therapeutics has historically been long, expensive, and fraught with uncertainty. Bringing a single new molecular entity from initial target identification to
regulatory approval typically spans ten to fifteen years and demands capital
investment exceeding two billion US dollars, with clinical attrition rates
consistently surpassing ninety percent across all therapeutic
areas~\cite{Paul2010, Pammolli2011}. The principal drivers of late-stage failure are insufficient efficacy and
unanticipated toxicity, both of which reflect an incomplete understanding of
molecular recognition at the target level during the earliest phases of
candidate selection~\cite{Paul2010}.
These figures collectively motivate the development of more rational,
computationally guided strategies for molecular design that can reduce the
time and cost burden at the hit-identification and lead-optimisation stages. For much of the twentieth century, medicinal chemistry and structure-based design dominated the search for bioactive compounds. While early advances in quantitative structure-activity relationship (QSAR) models and molecular docking partially mitigated costs, hypothesis generation remained largely manual, and the industry's productivity remained flat for decades~\cite{Lindsay1980}. However, the recent integration of deep generative models and self-supervised learning specifically masked representation learning has begun to reverse these trends. Recent analyses show that AI-native pipelines can compress the timeline from target discovery to Phase I clinical trials to under 24 months~\cite{Zhou2024}. Furthermore, AI-designed candidates have demonstrated Phase I success rates of $80\%$--$90\%$, significantly outperforming the historical industry average of approximately $40\%$~\cite{Jayatunga2024}. This shift is driven by the ability of generative architectures to navigate vast chemical spaces and predict complex absorption, distribution, metabolism, excretion, and toxicity (ADMET) profiles with unprecedented accuracy, moving beyond simple docking toward holistic, data-driven molecular evolution.


Early applications of AI in chemistry, exemplified by the \textit{DENDRAL} expert system, demonstrated that rule-based reasoning could assist chemists in interpreting mass spectra and predicting molecular structures~\cite{Lindsay1980}. However, it was not until the deep learning revolution of the 2010s that AI began to show genuine potential in molecular discovery. 

The computational representation of molecular structures is foundational to
both machine learning workflows and quantitative similarity-based drug
discovery.
The Simplified Molecular Input Line Entry System (SMILES) encodes the atomic
connectivity of a molecule as a linear string of characters, providing a
compact and machine-readable format for chemical structure
specification~\cite{Weininger1988} that has become the standard input for
cheminformatics and deep learning
pipelines~\cite{chithrananda2020chemberta, Ross2022}.
For similarity-based retrieval and diversity assessment,
extended-connectivity fingerprints (ECFPs), formally derived from the Morgan
algorithm~\cite{morgan1965}, encode the circular atomic neighbourhood around
each non-hydrogen atom up to a specified radius into fixed-length binary
vectors; the resulting representations capture local structural features that
are empirically correlated with biological activity~\cite{RogersHahn2010}.
The Tanimoto coefficient defined as the ratio of the number of bits set in
both fingerprints to the number of bits set in either fingerprint is the
most widely adopted metric for quantifying pairwise molecular similarity from
binary representations~\cite{Bajusz2015}.
Interpretation of specific Tanimoto values is, however, strongly dependent
on the molecular representation employed: Maggiora et al.\ critically
examined the widely propagated threshold of $\mathrm{T_c} \geq 0.85$ as a
universal indicator of shared bioactivity, demonstrating that this value
originated from a single study using proprietary fingerprints and cannot be
reliably generalised across different fingerprint types or compound activity
classes~\cite{maggiora2014}.
Observed Tanimoto distributions for confirmed active compound pairs vary
markedly by target, with median values spanning approximately 0.3 to 0.75
under MACCS fingerprints alone across ten exemplary activity
classes~\cite{maggiora2014}; accordingly, the coefficient is most robustly
interpreted on a relative scale for compound ranking and structural diversity
assessment rather than as an absolute activity predictor~\cite{maggiora2014}. In the present pipeline, Tanimoto similarity computed over Morgan fingerprints serves as the primary metric for quantifying the internal diversity of generated molecule sets, enabling principled comparison of the structural spread produced under interaction-aware versus random masking conditions

The application of deep generative models to molecular design represents a
conceptual shift from searching predefined compound libraries towards actively
constructing novel chemical matter with targeted
properties~\cite{GomezBombarelli2018}.
Variational autoencoders (VAEs) introduced by Kingma and
Welling~\cite{Kingma2013} were among the first deep generative architectures
adapted to molecular design; by encoding SMILES strings into a continuous
latent space, VAEs enable gradient-based traversal toward molecules with
optimised physicochemical properties~\cite{GomezBombarelli2018}.
Generative adversarial networks (GANs)~\cite{Goodfellow2014} provided an
alternative generative framework in which a generator network is trained
against a discriminator to produce structurally realistic molecular graphs
without requiring an explicit likelihood model.
Recurrent neural networks (RNNs) trained on SMILES corpora demonstrated early
on that sequence-based language models could autonomously generate drug-like
molecules with structural characteristics resembling known bioactives,
establishing SMILES generation as a viable \textit{de novo} design
paradigm~\cite{Segler2018}.
Subsequent integration of reinforcement learning with these generative
frameworks enabled multi-objective molecular optimisation, simultaneously
directing generation toward desired profiles of potency, selectivity,
solubility, and synthetic accessibility~\cite{Olivecrona2017}.


The transformer architecture, introduced for sequence transduction tasks by
Vaswani et al.~\cite{Vaswani2017} and extended to contextualised language
representation viaasked language modelling in BERT~\cite{Devlin2018},
established a new paradigm for learning deep bidirectional representations
from sequential data.
ChemBERTa adapted the RoBERTa pre-training framework to molecular SMILES
strings, employing either a byte-pair encoder or a SMILES-specific tokeniser
to construct a chemical vocabulary and training on randomly masked token
prediction over unlabelled compound datasets from
PubChem~\cite{chithrananda2020chemberta}.
This approach yields token-level molecular representations; the authors
hypothesised that, in learning to recover masked tokens, the model forms a
representational topology of chemical space that generalises to downstream
property prediction, and the resulting models demonstrated performance that
scales consistently with pretraining corpus size across MoleculeNet benchmarks
including blood--brain barrier permeability, HIV inhibition, and clinical
toxicity classification~\cite{chithrananda2020chemberta}.
Ross et al.\ subsequently demonstrated that scaling self-supervised SMILES
pretraining to a substantially larger corpus of 1.1~billion molecules from
PubChem and ZINC, using a transformer encoder with linear attention and
rotary positional embeddings, produces molecular representations that
outperform ChemBERTa and graph neural network baselines across a broad range
of MoleculeNet classification and regression benchmarks, including aqueous
solubility, lipophilicity, and quantum-chemical property prediction,
confirming that large-scale pretraining on SMILES captures rich structural
and physicochemical information without labelled activity data~\cite{Ross2022}.


In its standard formulation, however, the MLM objective selects positions for
masking uniformly at random (typically 15\% of tokens per input sequence) 
and trains the model to recover the original tokens conditioned on the unmasked
context~\cite{Devlin2018, chithrananda2020chemberta}. While this strategy successfully yields 
bidirectional, atom-level representations that generalise across molecular property 
prediction benchmarks~\cite{chithrananda2020chemberta, Ross2022}, its indifference 
to chemical semantics imposes a fundamental limitation: all atomic positions within a 
SMILES string are treated as equally informative targets, irrespective of whether they 
constitute the pharmacophoric core of a known binder or a chemically inert peripheral 
substituent. This limitation becomes particularly consequential when the MLM mechanism 
is repurposed not for representation learning but for conditional molecular 
reconstruction (the application at the heart of the present work).


Recognition that random masking fails to exploit the hierarchical chemical structure of 
molecules has motivated a wave of structured masking strategies. The MLM-FG framework 
demonstrated that replacing random token selection with masking of entire SMILES 
subsequences corresponding to predefined functional groups compels the model to 
reconstruct chemically coherent units rather than isolated atoms, yielding significant 
improvements across 11 molecular property prediction benchmarks relative to random and 
subsequence masking controls~\cite{peng2025mlmfg}. Complementary work introduced FARM, 
which combines MLM on SMILES with graph neural network encoders and contrastive 
learning, showing that functional-group-aware masking at the atom level achieves 
state-of-the-art performance on 10 of 12 MoleculeNet tasks by forcing the model to 
capture both local chemical environment and global molecular topology 
simultaneously~\cite{nguyen2024farm}. Systematic investigation of masking 
hyperparameters further revealed that molecular language models benefit from 
substantially higher masking ratios than those used in natural language processing, 
with optimal performance observed at approximately 35\% token masking rather than the 
conventional 15\%, a finding that reflects the denser information content of chemical 
tokens relative to natural language words~\cite{kruger2025molencoder}. Collectively, 
these studies establish that the identity and granularity of masked positions rather 
than just the masking ratio critically determines what structural abstractions a 
molecular language model learns to reconstruct.


A parallel line of work has sought to incorporate biological context into the generative 
process by conditioning molecular construction on pharmacophoric or pocket-level 
information. Zhu et al.~\cite{zhu2023pgmg} proposed PGMG, in which RDKit-identified 
chemical feature types hydrogen bond donors and acceptors, hydrophobic groups, and 
aromatic rings and their pairwise molecular-graph distances are encoded as a fully 
connected graph by a gated graph convolutional network and used to condition a 
variational transformer decoder for autoregressive SMILES generation, producing 
molecules with docking affinities comparable to known actives across fifteen targets 
including BRD4 without requiring explicit 3D protein structures or target-specific 
activity data. More recently, masked discrete 
diffusion frameworks such as GenMol have extended this paradigm by treating SMILES 
generation as a masked token infilling problem amenable to fragment-constrained and 
goal-directed tasks simultaneously, achieving state-of-the-art performance across 
multiple drug discovery scenarios without task-specific fine-tuning~\cite{lee2025genmol}. 
At the three-dimensional level, Lingo3DMol integrated a separately trained non-covalent 
interaction predictor into a language model for pocket-based 3D molecule generation, 
demonstrating that explicitly modelling hydrogen bonds and hydrophobic contacts during 
the generative step improves both pose quality and binding affinity of generated 
candidates~\cite{feng2024lingo3dmol}.



Bromodomain-containing protein 4 (BRD4), a member of the bromodomain and
extra-terminal domain (BET) family of epigenetic reader proteins, has emerged
as a compelling oncological target owing to its central role in the
transcriptional regulation of proto-oncogenes including \textit{MYC} and
\textit{BCL2}~\cite{filippakopoulos2010selective, ShiVakoc2014}.
BET bromodomains recognise and bind acetylated lysine residues on histone
tails vNIP Masking a conserved hydrophobic pocket, thereby recruiting transcriptional
machinery to active chromatin loci~\cite{filippakopoulos2010selective}.
The development of JQ1, a potent and selective thienodiazepine-based BET
bromodomain inhibitor, provided the first pharmacological demonstration that
competitive displacement of BRD4 from acetylated histones suppresses
oncogenic transcription and induces apoptosis in haematological
malignancies~\cite{filippakopoulos2010selective, Delmore2011}.
The well-defined acetyl-lysine binding pocket, the availability of numerous
high-resolution co-crystal structures, and the clinical relevance of BRD4
inhibition collectively make it an ideal model system for developing and
benchmarking structure-guided molecular generation
strategies~\cite{ShiVakoc2014}.

Molecular docking has served as a cornerstone of structure-based drug
discovery for several decades, providing a computational framework for
predicting the binding pose of a small molecule within a protein active site
and estimating the associated binding affinity~\cite{Kitchen2004}.
Docking algorithms sample the conformational and translational degrees of
freedom of a ligand within a defined receptor cavity and rank candidate poses
using empirical, force-field-based, or knowledge-based scoring
functions~\cite{Kitchen2004}. Virtual screening workflows built upon docking enable the prioritisation of
active compounds from libraries numbering in the millions, substantially
reducing the experimental burden at the hit-identification
stage~\cite{Shoichet2004}.
More recently, deep learning-based scoring functions have demonstrated
improved pose-prediction accuracy by training directly on large repositories
of protein--ligand structures; GNINA implements a convolutional neural network
scoring function that outputs both a pose quality score (CNNscore) and a
predicted binding affinity (CNNaffinity), and has shown competitive
performance against classical scoring functions on standard
benchmarks~\cite{mcnunn2021gnina, Dunn2024CACHE}.

The reliable identification and classification of non-covalent interactions at
protein--ligand interfaces is essential for understanding molecular recognition
and for guiding structure-based design.
The Protein--Ligand Interaction Profiler (PLIP) is an open-source tool that
automatically detects and characterises seven classes of non-covalent
interactions hydrogen bonds, hydrophobic contacts, salt bridges, water
bridges, $\pi$-stacking, $\pi$-cation interactions, and halogen bonds
directly from Protein Data Bank (PDB) coordinate files, without requiring
manual structure preparation~\cite{salentin2015plip}.
PLIP applies mostly geometry-based criteria with defined distance and angle
cutoffs derived from the structural biology literature to classify each
interaction type, and outputs per-atom interaction annotations that include
donor and acceptor identities, interatomic distances, and interaction
geometry~\cite{salentin2015plip}.
PLIP~2021 extended the original implementation with support for nucleic acid
targets, adjustable interaction detection thresholds, and flexible mode and
model selection, while retaining integration into high-throughput
computational pipelines its command-line interface and containerised
deployment~\cite{Adasme2021}.
The per-atom resolution of PLIP output specifying precisely which ligand
atoms participate in each interaction class and makes it directly applicable
to the construction of interaction-aware masking schemes for generative
molecular modelling. 




Returning to the generative modelling landscape, despite this convergence of structured masking and pharmacophore-conditioned generation, 
a precise mechanistic gap persists. Docking-based virtual screening, while capable of evaluating molecular fitness
post hoc, does not itself direct the generative process toward interaction-complementary
chemotypes~\cite{Kitchen2004}; nor do existing structured masking strategies. Existing structured masking strategies operate on 
chemical heuristics functional group membership, ring systems, or statistical 
co-occurrence patterns in SMILES corpora that are defined entirely in the absence of 
target-specific structural data \cite{peng2025mlmfg, nguyen2024farm, kruger2025molencoder}. 
Conversely, pocket-conditioned generative models typically require either explicit 3D 
protein structures \cite{peng2022pocket2mol, feng2024lingo3dmol}, pre-specified 
pharmacophore hypotheses derived from expert knowledge \cite{langer2010pharmacophore, 
zhu2023pgmg}, or substantial labelled activity data to train interaction predictors 
\cite{Stokes2020} requirements that limit their applicability at the hit-identification 
stage where structural information for novel targets may be sparse 
\cite{bender2021ai, Jumper2021}. No prior framework has directly coupled the per-atom 
non-covalent interaction annotations produced by automated protein--ligand interaction 
profilers \cite{salentin2015plip, Adasme2021} to a selective atom-level masking scheme 
within a pre-trained molecular language model \cite{chithrananda2020chemberta}, thereby 
converting co-crystallographic structural data into a zero-shot pharmacophoric prior 
without requiring task-specific fine-tuning or explicit 3D generation.

The present work closes this gap. By coupling PLIP-derived interaction 
annotations~\cite{salentin2015plip, Adasme2021} which specify precisely which ligand 
atoms participate in hydrogen bonds, water bridges, hydrophobic contacts, and 
$\pi$-stacking with the masked language modelling reconstruction mechanism of 
ChemBERTa~\cite{chithrananda2020chemberta}, we introduce an interaction-aware (NIP masking ) 
masking strategy in which biologically validated pharmacophoric positions are selectively 
occluded and the model is prompted to regenerate them conditioned on their full molecular 
context. This approach is conceptually complementary to functional-group-aware 
masking~\cite{peng2025mlmfg} but operates at a qualitatively different level of 
biological specificity: rather than masking chemically defined motifs uniformly across a 
compound library, NIP  masking targets the subset of atoms that are empirically engaged in 
productive non-covalent contacts within a given protein binding pocket, making the masked 
positions structurally and pharmacologically specific to the target of interest. The 
resulting pipeline GenPLIP is evaluated against random masking across five BRD4 
co-crystal ligand systems spanning distinct chemotypes and interaction profiles, 
benchmarked for chemical validity, internal diversity, ChEMBL pharmacological relevance, 
and GNINA docking performance, to quantify whether interaction-guided site selection 
constitutes a meaningful biological prior for SMILES-space molecular exploration.








\section{Methodology}
\subsection{Interactive Pipeline}

The masking pipeline is a comprehensive system for identifying protein-ligand interaction sites from PLIP (Protein-Ligand Interaction Profiler) data and masking corresponding atoms in molecular representations (SMILES/SELFIES). This enables machine learning models to learn context-aware molecular generation conditioned on protein binding pocket interactions. The pipeline architecture is structured like a waterfall model and consists of 14 different steps.

\subsubsection{Extraction of Ligand-Specific Interactions}
The first operational stage is the identification of interactions between a
chosen ligand and its surrounding protein environment. PLIP provides a
structured interaction report in which each interaction type is described
separately, allowing the ligand to be analyzed not as a generic chemical object
but as a target-specific binding partner.

For each interaction class $T$, for example, hydrogen bonds, hydrophobic
contacts, $\pi$-stacking, $\pi$-cation interactions, or salt bridges, PLIP
identifies a set of ligand atoms that participate in that interaction. We denote
this ligand-atom set by $I_T$. The value of this step is that it isolates
exactly which ligand atoms are being used by the protein to stabilize the
complex.

In order, this is the point where the ligand becomes
interaction-aware. Rather than asking which atoms exist in the
molecule, the pipeline asks which atoms are relevant to binding. That
distinction is essential because the later masking decisions are based on
interaction participation rather than molecular size or atom count alone.

% ─────────────────────────────────────────────
\subsubsection{Map PDB Serial Numbers for Specific Interactions}

Once PLIP has identified the interacting ligand atoms, the next challenge is
reconciling those atoms with the atom identities recorded in the PDB structure.
PDB files use serial numbers, whereas molecular toolkits often operate with
their own internal atom indices. In addition, a ligand may be reordered when
loaded into a cheminformatics toolkit, so a direct index-by-index comparison is
unreliable.

To address this issue, the pipeline uses a coordinate-based matching
strategy. Let $P = \{\vec{x}_{p,1}, \ldots, \vec{x}_{p,n}\}$ be the set of
3D coordinates for atoms extracted directly from the PDB ligand, and let
$R = \{\vec{x}_{r,1}, \ldots, \vec{x}_{r,n}\}$ be the coordinates of the same
atoms within the reconstructed ligand object. The mapping function
$f: \text{Serial}_{\text{PDB}} \to \text{Index}_{\text{RDKit}}$ identifies the
toolkit atom index that lies closest to each PDB atom in 3D space:

\begin{equation}
    f(i) = \arg\min_{j} \left\| \vec{x}_{p,i} - \vec{x}_{r,j} \right\|
    \label{eq:coord_mapping}
\end{equation}

\noindent
Here, $i$ indexes atoms by their PDB serial number, $j$ ranges over all atom
indices in the reconstructed RDKit ligand object, and
$\left\| \vec{x}_{p,i} - \vec{x}_{r,j} \right\|$ is the Euclidean distance
between the $i$-th PDB atom and the $j$-th RDKit atom in 3D space. For each
PDB atom $i$, the function $f(i)$ returns the RDKit index $j$ whose 3D
position is closest, thereby establishing a one-to-one correspondence between
The two numbering systems.

The mapping is accepted only if the coordinate discrepancy remains below a
strict tolerance threshold $\varepsilon$:

\begin{equation}
    \left\| \vec{x}_{p,i} - \vec{x}_{r,f(i)} \right\| < \varepsilon
    \label{eq:tolerance}
\end{equation}

\noindent
Equation~\ref{eq:tolerance} acts as a validation gate: after the nearest
neighbour $f(i)$ has been identified via Equation~\ref{eq:coord_mapping}, The 
residual distance between the matched pair must not exceed $\varepsilon$.
If the condition is violated for any atom, the match is rejected and flagged
as a mapping failure, preventing erroneous atom correspondences from
propagating silently through the pipeline. We use $\varepsilon = 0.15$~\AA,
a threshold intentionally tight so that the atom correspondence is not merely
approximate, but structurally trustworthy.

This step is the bridge between the PDB serial
numbering system and the ligand index system used later in the representation
stage.


% ─────────────────────────────────────────────
\subsubsection{Extraction of a Specific Ligand from a PDB File}

After the interacting atoms have been identified and aligned, the ligand itself
must be extracted as a standalone chemical entity from the protein-ligand
complex. This extraction is not just a file-splitting operation; it is the point
at which the ligand is separated from the protein environment while retaining
It's geometric and connectivity information.

The extracted ligand should preserve the atom ordering required for later
mapping, as well as the 3D coordinates needed for interaction reasoning. In the
context of a single protein-ligand complex, this ensures that the same ligand
can be traced across three different views: the crystallographic PDB record, the
PLIP interaction annotation and the molecular representation used for masking.

This stage is especially important when a PDB file contains multiple hetero
groups or when a ligand appears in a region with water molecules, cofactors, or
alternate conformations. This step aims to isolate the target ligand
cleanly enough that downstream atom mapping remains unambiguous.

% ─────────────────────────────────────────────
\subsubsection{Preservation of CONECT Records for Intra-Ligand Bonds}

Extracting the ligand as a set of atomic coordinates is a necessary but
insufficient condition for a complete structural representation. PDB files
encode intra-ligand bonding explicitly through \texttt{CONECT} records, which
specify which atoms are covalently connected. These records carry
information that cannot be reliably inferred from coordinates alone, including
ring membership, valence relationships, and the overall topology of the chemical
skeleton.

Discarding \texttt{CONECT} records during extraction reduces the ligand to an
unconnected point cloud. In that degenerate state, downstream operations that
depend on graph-level molecular structure bond order correction, aromatic ring
recognition, and atom-level masking either produce incorrect results or fail. Retaining these records, therefore, acts as a structural safeguard,
ensuring that the bond graph present in the original crystallographic entry is
faithfully carried forward into every subsequent stage of the pipeline.

This step supports chemical continuity, the
guarantee that the intramolecular connectivity of the extracted ligand remains
consistent with the deposited structure throughout the entire pipeline. This
consistency is a prerequisite for any representation, whether SMILES or SELFIES,
that is expected to encode a chemically valid molecule rather than an arbitrary
collection of atoms.


% ─────────────────────────────────────────────
\subsubsection{Conversion of 3D Ligand Structure into 1D Molecular String}

After the ligand has been isolated and its connectivity preserved, the next
stage is to convert the 3D structure into a 1D molecular representation. This
conversion is necessary because the masking process is ultimately performed on
strings rather than on coordinate files.

SMILES and SELFIES serve complementary roles in this workflow. SMILES offers a
compact, widely used textual encoding of molecular structure, while SELFIES
provides a representation that is robust to invalid strings and, therefore, useful
for constrained generative tasks. The conversion step must therefore preserve as 
much chemically meaningful information as possible while translating the ligand
into a tokenized sequence.

This stage is also where the difference between geometry and
representation becomes explicit. The 3D structure carries interaction
context, but the 1D string is what will be presented to the model. The pipeline
keeps these two views aligned through atom mapping and careful reconstruction
of bonding information.

% ─────────────────────────────────────────────
\subsubsection{Mapping of PDB Serial Numbers to PLIP-Generated Index}

Although PLIP interaction reports are derived from the full protein--ligand
complex, the masking target is usually a single ligand associated with a single
protein. This requires the interaction annotations to be reconciled not only
with PDB serial numbers, but also with the ligand-level indices used in the
final molecular representation.

This stage creates a consistent chain of identity across the structural
workflow: the interaction atom identified by PLIP is associated with its PDB
serial number, then matched to the reconstructed ligand atom index, and finally
associated with the corresponding token position in the textual molecule
representation. In methodological terms, this is the point where the interaction
Information becomes specific enough to support deterministic masking.

This single-protein mapping step is important because it prevents ambiguity
when multiple chains, alternate locations, or repeated residue numbering might
otherwise complicate the interpretation of atom-level annotations. It ensures
that the ligand being masked is the same ligand that was originally annotated
by PLIP in the context of the chosen protein target.

% ─────────────────────────────────────────────
\subsubsection{Fix Bond Order Using PLIP Templates}

Even when the ligand has been extracted correctly, bond orders may not always
be directly recoverable from the original coordinate file. To support chemically
valid downstream representations, the pipeline uses template-based bond
order correction informed by the known chemistry of the ligand and the
interaction-aware reconstruction process.

This step is to restore the correct distinction between single, double, triple, and aromatic bonds, where that information may
be incomplete or ambiguous after extraction. The correct bond order is essential
because it affects aromaticity perception, functional group recognition, and
The eventual validity of the SMILES or SELFIES output.

This step is best understood as a normalization
stage, the extracted ligand is not only isolated, but chemically regularized.
By applying PLIP-aware templates and chemically consistent reconstruction rules,
The pipeline avoids distortions that would otherwise propagate into the
representation-level masking step.

% ─────────────────────────────────────────────
\subsubsection{Generate SMILES/SELFIES with Atom Mapping Numbers}

Once the ligand has been reconstructed correctly, the atom mapping numbers are
introduced to establish a one-to-one correspondence between the internal atom
identities and the symbols that appear in the string representation. This is
crucial for SMILES in particular, because the textual syntax does not naturally
expose atom positions in a way that is directly usable for masking.

For SELFIES, the atom-token correspondence is more transparent because tokens
are explicit. A mapping function $g(i)$ connects the $i$-th ligand atom to the
position of its corresponding token in the SELFIES sequence. In SMILES, atom
mapping numbers serve the same purpose by tagging atoms so that their identity
can be recovered in the token stream.

The conceptual result is that the string representation becomes
atom-addressable. That means the pipeline can later substitute masked
positions without losing track of which original atom each symbol refers to.
This step preserves interpretability and ensures that the structural meaning of the mask remains traceable back to the 3D ligand.

% ─────────────────────────────────────────────
\subsubsection{Calculation of the Initial Atoms to Be Masked}

The next stage is to define which atoms should be hidden before any chemical
grouping rules are applied. The initial mask set depends on whether the pipeline
is operating in attractive or non-attractive mode.

\paragraph{Attractive Masking.}
In attractive masking, the goal is to hide atoms that directly participate in
protein-ligand interactions. If the relevant interaction sets are collected
across all selected interaction types $T \in \mathcal{T}$, then the attractive
mask set is defined as:

\begin{equation}
    M_{\text{attr}} = \bigcup_{T \in \mathcal{T}} \text{atoms in } S_T
    \label{eq:attractive_mask}
\end{equation}

\noindent
Here, $\mathcal{T}$ denotes the collection of all selected interaction types
(e.g.\ hydrogen bonds, hydrophobic contacts, $\pi$-stacking), and $S_T$ is the
set of ligand atoms participating in interaction type $T$. The union operator
ensures that an atom is included in $M_{\text{attr}}$ if it participates in
\emph{at least one} of the selected interaction types, so no interaction-bearing
atom is omitted regardless of how many interaction classes it belongs to.

\paragraph{Non-Attractive Masking.}
In non-attractive masking, the logic is reversed. Rather than masking the
interaction atoms, the pipeline masks everything else and preserves the
interaction-bearing core:

\begin{equation}
    M_{\text{non\text{-}attr}} = V \setminus M_{\text{attr}}
    \label{eq:nonattractive_mask}
\end{equation}

\noindent
Here, $V$ denotes the complete set of heavy atoms in the ligand, and the set
difference operator $\setminus$ removes all atoms already identified as
interaction-bearing. The result is that $M_{\text{non\text{-}attr}}$ contains
precisely those atoms that do \emph{not} participate in any selected interaction,
i.e.\ the chemical periphery surrounding the binding core. Taken together,
Equations~\ref{eq:attractive_mask} and~\ref{eq:nonattractive_mask} are
complementary and satisfy $M_{\text{attr}} \cup M_{\text{non\text{-}attr}} = V$.

This stage is important because it defines the pedagogical objective of the
masked example. In attractive mode, the model is asked to reconstruct
binding-critical fragments. In non-attractive mode, the model is asked to infer
the chemical surroundings of a preserved binding core.

% ─────────────────────────────────────────────
\subsubsection{Grouping of Atoms According to Interactions}

The initial atom list is rarely suitable for masking as is, because masking a
single atom can violate chemical structure or remove the context of a functional
group. For this reason, the pipeline expands atom-level interaction annotations
into chemically meaningful groups.

For each interaction type $T$, the set of interacting atoms $I_T$ is expanded
into a group collection through a chemical rule-based expansion function
$E(I_T, T)$. The outcome is a set of chemical groups:

\begin{equation}
    \mathcal{G}_T = \{ G_1, G_2, \ldots, G_k \}
    \label{eq:group_collection}
\end{equation}

\noindent
Here, $\mathcal{G}_T$ is the collection of all chemically expanded groups
derived from interaction type $T$, and each $G_i$ is a discrete block of atom
indices that must be treated as an indivisible unit during masking. The index
$k$ denotes the total number of such groups produced by the expansion function
$E(I_T, T)$ for that interaction type. Importantly, a single interacting atom
may expand into a group $G_i$ containing several atoms if the chemical context
demands it, for example, when that atom belongs to an aromatic ring or a
charged functional group.

The expansion function $E(I_T, T)$ is interaction-aware: its behaviour depends
not only on which atoms are involved, but also on the type of interaction
$T$ that identified them. Concretely:

\begin{itemize}
    \item \textbf{Aromatic interactions} ($\pi$-stacking, $\pi$-cation): The 
    interacting atom is expanded to encompass the full aromatic ring system to
    which it belongs, since masking a subset of an aromatic ring produces a
    chemically invalid fragment.

    \item \textbf{Charged interactions} (salt bridges, ionic contacts): The
    expansion includes the entire functional group carrying the formal charge,
    such as a carboxylate or ammonium moiety, because the charge is a property
    of the group rather than of any single atom.

    \item \textbf{Hydrophobic and other contacts}: In the absence of a broader
    structural rule, the interacting atom may remain as a singleton group
    $G_i = \{i\}$, unless it is part of a ring or functional group that the 
    higher priority rule would expand.
\end{itemize}

This stage is at which ring systems, functional groups, and local
interaction neighbourhoods are preserved as coherent masking units. The goal is
not merely to hide atoms, but to hide chemically meaningful fragments in
a way that remains interpretable to a molecular model.

% ─────────────────────────────────────────────
\subsubsection{Application of Chemical Rules for Rings and Functional Groups}

Masking individual atoms within a chemically coherent structural motif such as an aromatic ring or a carboxylate group may produce partial masks that are chemically uninformative or syntactically ambiguous in the molecular string. To avoid this, the pipeline applies a set of chemical rules that expand single-atom mask candidates to encompass their containing ring systems or functional groups. Ring membership is determined  RDKit's ring-perception algorithm, while functional group boundaries are identified using a curated SMARTS pattern library. Atoms added by rule expansion inherit the interaction group label of the atom that triggered the expansion.

% ─────────────────────────────────────────────
\subsubsection{Selection of Atoms According to Desired Percentage with Grouping Constraints}

The user specifies a desired masking fraction $P_T \in [0, 1]$ for each
interaction class $T$. The target number of atoms to mask is then computed as:

\begin{equation}
    D_T = \left\lceil P_T \times \sum_{G \in \mathcal{G}_T} |G| \right\rceil
    \label{eq:target_atoms}
\end{equation}

\noindent
Here, $P_T$ is the user-specific masking fraction for interaction type $T$,
$\mathcal{G}_T$ is the full collection of chemically expanded groups from
Step~10, and $|G|$ denotes the number of atoms in group $G$. The sum
$\sum_{G \in \mathcal{G}_T} |G|$ therefore gives the total number of atoms
available for masking across all groups of type $T$. The ceiling operator
$\lceil \cdot \rceil$ ensures that the target is always rounded up to the
nearest integer, so that the requested fraction is met rather than
underestimated when the product $P_T \times \sum |G|$ is non-integer.

Because the pipeline masks whole groups rather than isolated atoms, the selected
groups are accumulated greedily until the target $D_T$ is met or exceeded. This
means the exact requested percentage may not always be achieved, but the
chemical integrity of the selected fragments is preserved. To quantify the
actual result, the effective masking percentage is calculated as:

\begin{equation}
    P_{\text{eff},T} =
    \frac{\left| \bigcup_{G \in \mathcal{S}_T} G \right|}
         {\left| \bigcup_{G \in \mathcal{G}_T} G \right|}
    \label{eq:effective_percentage}
\end{equation}

\noindent
Here, $\mathcal{S}_T \subseteq \mathcal{G}_T$ is the subset of groups actually
selected for masking, and $\mathcal{G}_T$ is the full group collection. The
numerator counts the total number of distinct atoms covered by the selected
groups, while the denominator counts the total number of distinct atoms
available across all groups of type $T$. The union operators in both terms
ensure that atoms belonging to multiple groups are not double-counted.
Taken together, Equations~\ref{eq:target_atoms} and~\ref{eq:effective_percentage}
characterise the selection process: Equation~\ref{eq:target_atoms} sets the integer target, and Equation~\ref{eq:effective_percentage} measures how closely that target was achieved after group-level constraints were applied.

If $P_{\text{eff},T} > P_T$, the discrepancy typically reflects the fact that a
chemically large unit, such as an aromatic ring system or a charged functional group
had to be selected in its entirety to preserve structural consistency,
causing the masked atom count to overshoot the requested threshold.

% ─────────────────────────────────────────────

\subsubsection{Perform String-Level Masking}

After group selection, the actual string-level masking is performed on the
atom-mapped molecular representation produced in Step~8. The procedure differs
slightly depending on the chosen representation:

\begin{itemize}
    \item \textbf{SELFIES}: The atom-token mapping $g(i)$ is used to locate the
    token positions corresponding to the selected atoms, and each such position
    is replaced with the designated mask symbol. Because SELFIES tokens are explicit and non-overlapping, this substitution is unambiguous.

    \item \textbf{SMILES}: Atom mapping numbers identify the relevant atoms
    within the token stream. The corresponding token segments are substituted
    with the mask token, and residual mapping labels are removed from the
    output to produce a clean masked string. Care is taken to preserve the
    surrounding syntactic context so that the resulting string remains
    well-formed.
\end{itemize}

\noindent
The result in both cases is a masked sequence that still encodes the original
molecular context, but with the desired structural information hidden, ready
for use as a conditioned input to a generative model.

% ─────────────────────────────────────────────

\subsubsection{Generate a Detailed Report}

The final reporting step produces a structured summary of the entire masking
process. The report includes:

\begin{itemize}
    \item The interaction classes considered and their associated atom groups,
    \item The requested masking fraction $P_T$ and the achieved effective
    percentage $P_{\text{eff},T}$ for each interaction type $T$,
    \item The molecular representation type used (SMILES or SELFIES) and the
    final masked string,
    \item Any mapping ambiguities, structural exclusions, or tolerance
    violations encountered during the pipeline run.
\end{itemize}

\noindent
This report makes the masking process transparent and reproducible. When masked examples are used for model training, benchmarking, or qualitative analysis, the report provides the audit trail needed to reconstruct exactly which atoms were masked, why they were selected, and how
closely the outcome matched the intended masking fraction.

% ─────────────────────────────────────────────

\pagebreak
\subsection{ChemBERTa inference }
Masked SMILES strings were provided to a pretrained ChemBERTa model in masked language modeling mode. The model predicted replacements for [MASK] tokens, producing candidate completed SMILES strings. To evaluate task adaptation without full retraining, partial fine-tuning was additionally performed by unfreezing only the final transformer layers while keeping the remaining layers frozen, and fine-tuning on the masked dataset. Model behavior was compared between (i) inference with the fully frozen pretrained model and (ii) inference after partial fine-tuning. Baseline runs using unmasked SMILES were included to characterize model behavior prior to interaction-guided masking. 
Masked molecular sequence completion was performed using a pretrained ChemBERTa masked language model operating on SMILES representations. The input to the model could be provided either as SMILES or as SELFIES containing one or more mask tokens. To ensure compatibility with ChemBERTa, all inputs were first normalized to masked SMILES format using a preprocessing step that converts SELFIES to SMILES while preserving mask positions.
For a given masked SMILES input, tokenization was performed using the corresponding ChemBERTa tokenizer, yielding token IDs and attention masks. A forward pass through the model produced position-wise logits over the vocabulary for each token in the sequence.
Positions corresponding to the model’s mask token were identified by matching against the tokenizers mask token id.
For each masked position, the output logits were converted to probabilities using a softmax operation. The top K most probable token candidates were retained, and their probabilities were renormalized to define a categorical distribution. These per-mask distributions were precomputed to enable efficient sampling.
Candidate molecular sequences were generated by independently sampling a replacement token for each masked position from its corresponding top K distribution. This process was repeated multiple times to generate a diverse set of candidate completions. Each sampled token sequence was then decoded back into a SMILES string using the tokenizer’s decoding function.


\subsection{Post-processing and validation}
Generated SMILES candidates were filtered and validated using RDKit and Open Babel. Validity filtering removed unparsable SMILES and structures failing basic chemical constraints. Feasibility checks included sanitization and structural consistency verification. Only molecules that could be successfully parsed into valid RDKit molecular graphs were retained. Duplicate valid SMILES were removed, and the final set of unique, chemically valid molecules was optionally written to disk for downstream analysis. The number of valid and invalid generations was tracked to assess sampling efficiency.
The resulting generated molecules were additionally assessed by graphical inspection against known/ground-truth ligand structures (e.g., visual comparison of 2D depictions and, where applicable, structural overlays) to verify preservation of key substructures and overall chemical plausibility. Recovery of identical or closely similar structures was used as supporting evidence that the interaction-guided masking and generation pipeline produced chemically plausible outputs consistent with the reference ligand chemistry.
 
\subsection{Molecular fingerprinting and chemical similarity}
 
\paragraph{Morgan circular fingerprints.}
All pairwise and nearest-neighbour chemical similarity calculations
were performed using Morgan circular fingerprints~\cite{morgan1965},
as implemented in RDKit~\cite{rdkit}.
For each molecule represented as a SMILES string, a circular
fingerprint of radius $r = 2$ and bit-vector length $L = 2048$ was
computed.
The Morgan algorithm assigns an initial invariant to each heavy
atom $a_i$ based on its atomic number, degree, number of hydrogens,
charge, and isotope label.
At each iteration $k$ (up to the specified radius), the invariant
of atom $a_i$ is updated by hashing its current invariant together
with the ordered invariants of all neighbouring atoms within $k$
bonds:
 
\begin{equation}
    \text{inv}^{(k)}(a_i) =
        \text{hash}\!\left(
            \text{inv}^{(k-1)}(a_i),\;
            \left\{
                \text{inv}^{(k-1)}(a_j) : a_j \in \mathcal{N}(a_i)
            \right\}
        \right),
    \label{eq:morgan_update}
\end{equation}
 
\noindent where $\mathcal{N}(a_i)$ denotes the set of atoms
directly bonded to $a_i$.
After $r = 2$ iterations, the set of all atom-environment
identifiers $\{\text{inv}^{(k)}(a_i)\}$ for $k = 0, 1, 2$ is
folded into a binary bit-vector $\mathbf{f} \in \{0,1\}^{2048}$
by mapping each identifier to a bit position  modular hashing:
 
\begin{equation}
    f_b = 1 \iff \exists\; (i, k) \;\text{s.t.}\;
          \text{inv}^{(k)}(a_i) \bmod L = b,
    \label{eq:morgan_fold}
\end{equation}
 
\noindent where $b \in \{0, \ldots, L-1\}$ is a bit index.
The resulting fingerprint encodes the presence or absence of all
circular substructures up to two bonds from every heavy atom,
capturing both local chemical environment and connectivity topology.
 
\paragraph{Tanimoto similarity.}
Chemical similarity between any two molecules $m_i$ and $m_j$
with binary fingerprint vectors $\mathbf{f}_i$ and $\mathbf{f}_j$
was computed using the Tanimoto (Jaccard) coefficient:
 
\begin{equation}
    T(\mathbf{f}_i, \mathbf{f}_j)
    = \frac{|\mathbf{f}_i \cap \mathbf{f}_j|}
           {|\mathbf{f}_i \cup \mathbf{f}_j|}
    = \frac{\mathbf{f}_i \cdot \mathbf{f}_j}
           {|\mathbf{f}_i|^2 + |\mathbf{f}_j|^2
            - \mathbf{f}_i \cdot \mathbf{f}_j},
    \label{eq:tanimoto}
\end{equation}
 
\noindent where $\mathbf{f}_i \cdot \mathbf{f}_j = \sum_b f_{i,b}
\cdot f_{j,b}$ is the number of bits set in both fingerprints
(intersection count), and $|\mathbf{f}|^2 = \sum_b f_b$ denotes
the number of bits set in fingerprint $\mathbf{f}$ (popcount).
$T$ ranges from 0 (no shared substructures) to 1 (identical
fingerprints).
A threshold of $T \geq 0.40$ was applied as the standard
scaffold-similarity boundary in medicinal chemistry~\cite{maggiora2014},
below which two compounds are considered to belong to different
scaffold classes.
 
\paragraph{Pairwise Tanimoto analysis .}
For each generated library $\mathcal{G}$ produced under a given
masking strategy and masking depth, all $\binom{|\mathcal{G}|}{2}$
pairwise Tanimoto scores were computed to characterise the internal
chemical diversity of the library:
 
\begin{equation}
    \bar{T}_{\mathrm{pair}}(\mathcal{G})
    = \frac{2}{|\mathcal{G}|(|\mathcal{G}|-1)}
      \sum_{i < j} T(\mathbf{f}_i, \mathbf{f}_j).
    \label{eq:mean_pairwise}
\end{equation}
 
\noindent A lower $\bar{T}_{\mathrm{pair}}$ indicates higher
internal structural diversity.
Scores were aggregated across all masking depths $k = 1, \ldots, N$
to produce the distribution shown in
Fig.~\ref{fig:pairwise_tanimoto}.
Interaction-aware (NIP Masking) and random masking strategies were evaluated
independently, and the resulting distributions were compared to
assess whether PLIP-guided site selection affects the internal
chemical diversity of the generated libraries.
 
\subsection{ChEMBL nearest-neighbour analysis }
 
\paragraph{Reference database.}
All generated molecules were compared against the ChEMBL
database~\cite{gaulton2012chembl} (release 33, $\approx 2.4 \times
10^6$ compounds with reported bioactivity data).
The full ChEMBL compound set was pre-processed by standardising
SMILES with RDKit's \texttt{MolStandardize} module, removing
salts and mixtures, and retaining only compounds with valid
two-dimensional structures, yielding a reference fingerprint
matrix $\mathbf{F}_{\mathrm{ChEMBL}} \in \{0,1\}^{N_{\mathrm{ref}}
\times L}$.
 
\paragraph{Nearest-neighbour retrieval.}
For each generated molecule $m_i \in \mathcal{G}$, the
nearest neighbour in ChEMBL was defined as the compound
$m^{*}_i$ that maximises the Tanimoto similarity:
 
\begin{equation}
    m^{*}_i = \operatorname*{arg\,max}_{c \;\in\; \mathcal{C}}
              \; T(\mathbf{f}_i, \mathbf{f}_c),
    \label{eq:nn_retrieval}
\end{equation}
 
\noindent where $\mathcal{C}$ denotes the set of all ChEMBL
compounds.
Only molecules for which $T(\mathbf{f}_i, \mathbf{f}_{m^{*}_i})
\geq 0.40$ were retained for further analysis; molecules with no
ChEMBL neighbour above this threshold were classified as having no
known bioactive analogue.
 
\paragraph{Nearest-neighbour similarity histogram.}
The distribution of nearest-neighbour similarities across a
generated library was summarised as a histogram of
$\{T(\mathbf{f}_i, \mathbf{f}_{m^{*}_i})\}_{m_i \in \mathcal{G}}$
(Fig.~\ref{fig:nn_similarity_hist}).
This distribution characterises how close the generated library is
to known ChEMBL chemical space: molecules in the range
$T \in [0.40, 0.60)$ are classified as scaffold analogues with
low similarity (novel relative to known compounds), those in
$T \in [0.60, 0.80)$ as medium-similarity analogues, and those
with $T \geq 0.80$ as high-similarity analogues (near-duplicates
of known compounds).
 
\paragraph{Nearest-neighbour frequency chart.}
To identify which specific ChEMBL compounds are most frequently
matched by the generated library, the frequency of each unique
nearest neighbour $c \in \mathcal{C}$ was computed as:
 
\begin{equation}
    \text{freq}(c)
    = \left|\left\{
        m_i \in \mathcal{G} \;:\;
        m^{*}_i = c \;\text{ and }\;
        T(\mathbf{f}_i, \mathbf{f}_c) \geq 0.40
      \right\}\right|.
    \label{eq:nn_frequency}
\end{equation}
 
\noindent The top-ranked ChEMBL compounds by $\text{freq}(c)$ were
plotted as stacked bar charts (Fig.~\ref{fig:nn_frequency}), with
each bar partitioned into Low ($0.40 \leq T < 0.60$), Medium
($0.60 \leq T < 0.80$), and High ($T \geq 0.80$) similarity
tiers, and annotated with the median Tanimoto similarity over
all mapped generated molecules:
 
\begin{equation}
    \widetilde{T}(c)
    = \operatorname{median}
      \left\{
          T(\mathbf{f}_i, \mathbf{f}_c) :
          m^{*}_i = c,\;
          T(\mathbf{f}_i, \mathbf{f}_c) \geq 0.40
      \right\}.
    \label{eq:nn_median}
\end{equation}
 
\noindent NIP Masking and random masking libraries were analysed separately
and compared side by side to assess whether the masking strategy
influences the identity and similarity tier of the ChEMBL compounds
most frequently recovered as nearest neighbours.
 
\subsection{Molecular docking }
 
\paragraph{Receptor and ligand preparation.}
Receptor structures were prepared from the co-crystal PDB files
used in Stage~1.
All ATOM records were extracted to form the receptor coordinate
file, while the HETATM records matching the original ligand
residue name, chain identifier, and residue sequence number were
isolated as the autobox reference ligand.
Generated molecules were converted from canonical SMILES to
three-dimensional conformers using RDKit's ETKDGv3 embedding
algorithm~\cite{riniker2015}, followed by energy minimisation
with the Universal Force Field (UFF)~\cite{rappe1992uff}.
Molecules for which three-dimensional embedding failed were
discarded with a warning and excluded from all downstream
analyses.
 
\paragraph{Docking protocol.}
All generated molecules — from both interaction-aware (NIP Masking) and
random masking strategies, deduplicated by canonical SMILES —
were docked against the BRD4 binding pocket using GNINA~v1.0.3
(GPU-accelerated)~\cite{mcnunn2021gnina}.
The docking box was defined automatically (\texttt{--autobox\_ligand})
using the co-crystallised ligand as the reference, with default
GNINA autobox padding.
Up to nine binding poses were sampled per molecule.
Three scoring terms were extracted from GNINA's output SDF
for each pose:
 
\begin{itemize}
    \item \textbf{CNNscore} — a convolutional neural network (CNN)-predicted
          binding probability $p_{\mathrm{bind}} \in [0, 1]$, where
          higher values indicate a more plausible binding pose;
    \item \textbf{CNNaffinity} — a CNN-predicted binding affinity
          expressed as $-\!\log K_d$ (in equivalent units), where
          higher values indicate stronger predicted binding;
    \item \textbf{minimizedAffinity} (Vinardo score) — a physics-based
          empirical docking score in kcal/mol~\cite{quiroga2016vinardo},
          where more negative values indicate stronger predicted binding.
\end{itemize}
 
\noindent Only the top-ranked pose (pose~1) per molecule was retained
for all downstream analyses (Stage~7 and Stage~8).
Previously docked molecules were identified by the presence of a
non-empty, RDKit-parseable \texttt{docked\_poses.sdf} file and
served from cache without re-docking.
Docking results for all molecules across all ligand groups were
aggregated into a single \texttt{docking\_summary.csv} file
containing one row per (ligand group, molecule index, pose) with
columns for SMILES, CNNscore, CNNaffinity, and minimizedAffinity.
 
\paragraph{Preliminary ranking.}
For visual inspection, docked molecules were ranked within each
ligand group by CNNaffinity (primary key, descending) with
minimizedAffinity as tiebreaker (ascending), and a ranked PDF
table was produced.
 
\subsection{Top-docked molecule report }
 
For each BRD4 ligand group, the five highest-scoring generated
molecules (top-5 by CNNaffinity, using pose~1 only) were
identified and compiled into a structured PDF report.
Each page of the report showed: (i) a molecule image row
displaying the original co-crystal ligand alongside the
top-5 generated molecules, with the original ligand
cell bordered in green if it was reproduced among the top-5
and orange otherwise; and (ii) a scores table reporting the
full SMILES, CNNaffinity, CNNscore, and Vinardo score for each
top-5 molecule.
Standalone PNG images were also saved per ligand group for
inspection without opening the full PDF.
 
\subsection{Novelty--potency analysis}
 
\paragraph{Tanimoto similarity to the original ligand.}
For each docked generated molecule, the Tanimoto similarity
to the corresponding co-crystal reference ligand was computed
using Morgan fingerprints (radius $r = 2$, $L = 2048$ bits)
as described in the Methods section on molecular fingerprinting.
This per-molecule similarity value serves as the novelty axis
in the scatter plot visualisation: molecules with lower Tanimoto
similarity to the reference ligand are considered structurally
more novel.
 
\paragraph{Composite docking score.}
A composite score was computed for each molecule by taking the
product of CNNaffinity and CNNscore:
 
\begin{equation}
    S_{\mathrm{composite}}(m)
    = \text{CNNaffinity}(m) \times \text{CNNscore}(m),
    \label{eq:composite}
\end{equation}
 
\noindent which simultaneously rewards predicted binding affinity
and binding pose plausibility.
Both factors are higher-is-better, so their product is also
higher-is-better and is bounded below by zero.
$S_{\mathrm{composite}}$ was used as the primary potency axis
in scatter visualisations and as the objective function in
Pareto front analysis.
 
\paragraph{Hierarchical quality filter.}
Prior to virtual screening prioritisation, a two-step
hierarchical quality filter was applied to remove physically
implausible poses:
 
\begin{enumerate}
    \item Retain molecules with $\text{CNNscore} > \tau_1$
          (removes misplaced binding poses); default
          $\tau_1 = 0.70$.
    \item From the survivors, retain molecules with
          $\text{minimizedAffinity} < \tau_2$ kcal/mol
          (enforces a physical energy threshold); default
          $\tau_2 = -7.0$ kcal/mol.
    \item Sort the doubly-filtered set by CNNaffinity
          (descending) to produce the final ranked hit list.
\end{enumerate}
 
\noindent Formally, the filtered set $\mathcal{F}$ is:
 
\begin{equation}
    \mathcal{F} = \bigl\{
        m \in \mathcal{G} \;:\;
        \text{CNNscore}(m) > \tau_1 \;\wedge\;
        \text{minimizedAffinity}(m) < \tau_2
    \bigr\},
    \label{eq:hierarchical_filter}
\end{equation}
 
\noindent with $\mathcal{F}$ subsequently sorted by
$\text{CNNaffinity}(m)$ in descending order.
These thresholds are consistent with previously published
GNINA-based virtual screening workflows~\cite{mcnunn2021gnina}.
 
\paragraph{Pareto front analysis.}
To identify molecules that optimally balance structural novelty
and predicted binding potency, a Pareto front was computed over
two objectives:
 
\begin{align}
    \text{Objective 1 (potency):}  &\quad
        S_{\mathrm{composite}}(m) \to \text{maximise},
        \label{eq:obj1} \\
    \text{Objective 2 (novelty):}  &\quad
        \eta(m) = 1 - T\!\left(\mathbf{f}_m,\,
        \mathbf{f}_{\mathrm{ref}}\right) \to \text{maximise},
        \label{eq:obj2}
\end{align}
 
\noindent where $T(\mathbf{f}_m, \mathbf{f}_{\mathrm{ref}})$
is the Tanimoto similarity of the generated molecule $m$ to
the reference ligand (Eq.~\ref{eq:tanimoto}), and
$\eta(m) \in [0, 1]$ is the novelty score.
A molecule $m$ is defined as \emph{Pareto-optimal}
(non-dominated) if there exists no other molecule $m'$ in the
docked set $\mathcal{G}$ such that $m'$ is simultaneously
at least as good as $m$ on both objectives and strictly
better on at least one:
 
\begin{equation}
    m \in \mathcal{P} \iff
    \nexists\; m' \in \mathcal{G} \setminus \{m\} \;:\;
    \left[
        S_{\mathrm{composite}}(m') \geq S_{\mathrm{composite}}(m)
        \;\wedge\;
        \eta(m') \geq \eta(m)
    \right]
    \;\wedge\;
    \left[
        S_{\mathrm{composite}}(m') > S_{\mathrm{composite}}(m)
        \;\vee\;
        \eta(m') > \eta(m)
    \right].
    \label{eq:pareto}
\end{equation}
 
\noindent The Pareto front $\mathcal{P}$ was computed in
$O(|\mathcal{G}|^2)$ time by exhaustive pairwise dominance
checking, which is tractable for the library sizes generated
in this study ($|\mathcal{G}| \leq 10^3$ per ligand group).
The top-10 molecules from $\mathcal{P}$, sorted by
$S_{\mathrm{composite}}$ in descending order, were compiled
into a PDF report containing the 2D structure, SMILES,
CNNaffinity, CNNscore, Vinardo score, Tanimoto similarity
to the reference ligand, and composite score for each
candidate.
\subsubsection*{Butina scaffold-cluster
                representative selection}
 
For completeness we also implemented Butina scaffold-cluster which ensures scaffold-level structural diversity by first
clustering all generated molecules and then selecting one
representative per cluster.
Molecules were clustered using the Butina
algorithm~\cite{butina1999}, which partitions a compound set into
clusters based on an upper-bound Tanimoto distance threshold
$\varepsilon$ (default $\varepsilon = 0.40$).
 
The Butina algorithm proceeds as follows.
A pairwise Tanimoto distance matrix is computed for all $n$
generated molecules:
 
\begin{equation}
    d_{ij} = 1 - T(\mathbf{f}_i, \mathbf{f}_j),
    \qquad d_{ij} \in [0, 1].
    \label{eq:tanimoto_distance}
\end{equation}
 
\noindent Molecules are sorted by the number of neighbours within
distance $\varepsilon$ (i.e.\ by chemical connectivity), and the
molecule with the most neighbours is chosen as the first cluster
centroid.
All molecules within distance $\varepsilon$ of the centroid form
the first cluster and are removed from further consideration.
The process is repeated iteratively until all molecules are
assigned to a cluster.
This produces a partition $\{K_1, K_2, \ldots, K_C\}$ of the
generated library into $C$ clusters, where all molecules within
a cluster have pairwise Tanimoto similarity $\geq 1 - \varepsilon$
to the centroid.
 
From each cluster $K_c$, the molecule with the highest composite
score is selected as the representative:
 
\begin{equation}
    m^{*}_c
    = \operatorname*{arg\,max}_{m \in K_c}
      S_{\mathrm{composite}}(m).
    \label{eq:cluster_rep}
\end{equation}
 
\noindent The $C$ representatives are then sorted by composite
score in descending order and the top~10 are returned:
 
\begin{equation}
    \mathcal{C}_5
    = \operatorname{top\text{-}10}_{S_{\mathrm{composite}}\downarrow}
      \bigl\{ m^{*}_c \bigr\}_{c=1}^{C}.
    \label{eq:idea5}
\end{equation}
 
\noindent This strategy guarantees that no two candidates in the
top-10 belong to the same structural scaffold cluster, providing
the broadest scaffold-level chemical diversity of the four
strategies.
The number of clusters $C$ depends on both $\varepsilon$ and the
chemical diversity of the generated library; at $\varepsilon = 0.40$,
it is consistent with the Bemis--Murcko scaffold-equivalence
threshold used in the novelty analysis~\cite{bemis1996scaffolds}.
\paragraph{Scatter plot visualisation.}
Novelty--potency scatter plots were produced per ligand group
and as a combined cross-ligand panel, with the Tanimoto
similarity to the reference ligand on the $x$-axis and
$S_{\mathrm{composite}}$ on the $y$-axis.
A vertical dashed line at $T = 0.40$ marks the standard
scaffold-novelty threshold~\cite{maggiora2014}; molecules
to the left of this line are classified as novel-scaffold
candidates.
Pareto-optimal molecules are highlighted on each plot
to guide visual identification of the best novelty--potency
trade-off candidates.

\subsection{Property-guided supervised fine-tuning of ChemBERTa}
\label{sec:stage10}

The pipeline described so far uses ChemBERTa exactly as pretrained: masks are
placed by the interaction analysis and filled by the frozen model. This final
methodological stage asks whether the model itself can be improved, by
fine-tuning it so that the completions it proposes are more often valid,
drug-like, synthesisable and distinct from the ligand they were derived from.

\subsubsection{Rationale: selection in place of a score-function estimator}

The obstruction to optimising such a criterion directly is that it is not
differentiable. Validity, quantitative estimate of drug-likeness
(QED)~\cite{bickerton2012qed}, synthetic accessibility~\cite{ertl2009sascore},
Tanimoto novelty~\cite{RogersHahn2010} and the PAINS~\cite{baell2010pains} and
Brenk~\cite{brenk2008alerts} structural-alert filters are computed by a
cheminformatics parser acting on a \emph{decoded string}; the number RDKit
returns carries no gradient back to the network that produced the tokens.

That obstruction is a property of the \emph{scoring} function, not of the
\emph{training} objective, and it can therefore be circumvented without any
score-function (policy-gradient) estimator at all: sample several completions of
the same masked molecule, let the score decide which one is best, and fit the
model to that completion by ordinary maximum likelihood. The chemistry then
enters only through a discrete $\arg\min$, and every quantity that reaches
\texttt{backward()} is a log-probability produced by the network itself. This is
expert iteration~\cite{anthony2017expert} in its rejection-sampling form, also
known as self-taught or reward-ranked
fine-tuning~\cite{zelikman2022star,dong2023raft}; in molecular design it is the
hill-climbing fine-tuning used to build focused libraries from a pretrained
SMILES generator~\cite{Segler2018,brown2019guacamol}.

Two properties distinguish it from a policy gradient. First, the direction of
the update is bounded: a REINFORCE-style loss contains a term
$-(\mathrm{score}-b)\log P_{\theta}$, which for a below-baseline rollout is a
\emph{positive} multiple of $\log P_{\theta}$ and so drives the model's own
samples towards probability zero --- an objective unbounded below and a known
source of degenerate optima. Cross-entropy is bounded below by zero and only
ever raises the probability of tokens observed to produce a good molecule.
Second, the achievable optimum is capped by the sampler: selection can only
imitate the best member of a $K$-sample pool and cannot reach a completion the
model would not already propose. This stage accordingly trades reach for
stability.

\subsubsection{Composite selection loss}

Each candidate $\hat{s}$ was scored against its pre-mask parent $s$ by a
composite \emph{loss}, lower being better,
\begin{equation}
\ell(\hat{s}, s) \;=\;
\begin{cases}
\;w_{\mathrm{QED}}\bigl(1-\mathrm{QED}(\hat{s})\bigr)
\;+\; w_{\mathrm{SA}}\,\mathrm{clip}_{[0,1]}\!\left(\dfrac{\mathrm{SA}(\hat{s})-1}{9}\right)
\;+\; w_{\mathrm{nov}}\,T_{\mathrm{c}}(s,\hat{s})
\;+\; w_{\mathrm{alert}}\,\mathrm{alert}(\hat{s}),
& \hat{s}\ \text{valid},\\[2.2ex]
\;S \;+\; w_{\mathrm{valid}},
& \hat{s}\ \text{invalid},
\end{cases}
\label{eq:stage10-loss}
\end{equation}
where $\mathrm{QED}\in[0,1]$ is the drug-likeness estimate, $\mathrm{SA}\in[1,10]$
the synthetic-accessibility score, $T_{\mathrm{c}}(s,\hat{s})$ the Tanimoto
similarity between Morgan fingerprints of parent and candidate,
$\mathrm{alert}(\hat{s})\in\{0,1\}$ indicates a PAINS or Brenk match, and
$S \equiv w_{\mathrm{QED}} + w_{\mathrm{SA}} + w_{\mathrm{nov}} + w_{\mathrm{alert}}$.
The synthetic-accessibility term is normalised over the score's actual range as
$(\mathrm{SA}-1)/9$, so that the easiest-to-synthesise molecule is charged
exactly zero. The novelty term penalises \emph{similarity} directly, which makes
reconstruction of the parent ($T_{\mathrm{c}}=1$) the worst attainable value of
that term. Any component that RDKit parses but cannot evaluate is charged its
full weight, the pessimistic reading being the safe one for a criterion whose
only role is to rank candidates. Weights were $(w_{\mathrm{QED}},
w_{\mathrm{SA}}, w_{\mathrm{nov}}, w_{\mathrm{alert}}) = (0.30, 0.20, 0.30,
0.20)$, summing to $S=1$, and $w_{\mathrm{valid}}=1$.

Validity enters Eq.~\eqref{eq:stage10-loss} not as one weighted term among the
others but as a separating margin. Because the four property terms are
individually bounded by their weights,
\begin{equation}
\sup_{\hat{s}\ \mathrm{valid}} \ell(\hat{s},s) \;=\; S
\;<\; S + w_{\mathrm{valid}} \;=\; \ell(\hat{s},s)\big|_{\hat{s}\ \mathrm{invalid}}
\qquad \text{for any } w_{\mathrm{valid}}>0 ,
\label{eq:stage10-ordering}
\end{equation}
so \emph{every} valid candidate outranks \emph{every} invalid one, irrespective
of its drug-likeness, accessibility, novelty or alert status. Validity is thus a
lexicographic gate on the selection rather than a quantity traded off against
the others, and $w_{\mathrm{valid}}$ sets only the size of the margin, never the
ordering. Encoding the same priority by a large multiplicative penalty would
instead have made the selection criterion numerically incomparable with the
cross-entropy it feeds, which is of order $1$--$10$ nats.

\subsubsection{Candidate pool and target selection}

For each masked/parent pair $(\tilde{s}, s)$ a \emph{single} forward pass
produced the logits $z_i \in \mathbb{R}^{|V|}$ at every masked position $i \in
\mathcal{M}$. That one pass serves both roles in the algorithm: it defines the
proposal distribution the candidates are drawn from, and it supplies the logits
against which the cross-entropy is later computed, so no second forward pass is
required. $K=16$ completions were drawn independently from the same
temperature-scaled, top-$k$-truncated distribution used for inference,
\begin{equation}
\hat{t}^{(k)}_i \;\sim\; q_{\theta}\bigl(\cdot \mid c(\tilde{s})\bigr)
\;\propto\; \mathrm{softmax}\bigl(z_i / T\bigr)\Big|_{\mathrm{top}\text{-}k},
\qquad k = 1,\dots,K, \quad i \in \mathcal{M},
\label{eq:stage10-sampling}
\end{equation}
with $T=1.2$ and $k=20$, and each was decoded to a candidate SMILES
$\hat{s}^{(k)}$. Sampling was performed under \texttt{detach}: the draw is a
\emph{search} over possible targets and carries no gradient, which is precisely
what distinguishes this estimator from REINFORCE, where the sampled token's
log-probability is retained with its computation graph intact. The pool was
scored by Eq.~\eqref{eq:stage10-loss} and the winner taken as
\begin{equation}
k^{\star} \;=\; \arg\min_{k \in \{1,\dots,K\}} \ell\bigl(\hat{s}^{(k)}, s\bigr),
\qquad
\mathcal{V} \;=\; \mathbb{1}\bigl[\hat{s}^{(k^{\star})}\ \text{is valid}\bigr].
\label{eq:stage10-argmin}
\end{equation}
By Eq.~\eqref{eq:stage10-ordering}, $\mathcal{V}=1$ if and only if at least one
of the $K$ candidates parses.

\subsubsection{Parent fallback and its weighting}

Measured on these data at a masking rate of $p=15\%$, approximately $70\%$ of
molecules yield \emph{no} valid completion in $K=16$ draws, and $K=8$ gives an
indistinguishable rate --- validity is bimodal and molecule-intrinsic here, not
a matter of drawing enough samples. Those molecules offer nothing worth
imitating, so the target reverts to the parent's own tokens, which are available
for every pair and valid by construction:
\begin{equation}
\mathbf{t}^{\star} \;=\;
\begin{cases}
\bigl(\hat{t}^{(k^{\star})}_i\bigr)_{i \in \mathcal{M}}, & \mathcal{V}=1,\\[1ex]
\bigl(t^{\,s}_i\bigr)_{i \in \mathcal{M}}, & \mathcal{V}=0,
\end{cases}
\qquad
\omega \;=\;
\begin{cases}
1, & \mathcal{V}=1,\\[1ex]
\lambda_{\mathrm{fb}}, & \mathcal{V}=0,
\end{cases}
\label{eq:stage10-target}
\end{equation}
where $t^{\,s}_i$ is the token of the canonical parent at position $i$.
Recovering $t^{\,s}_i$ requires no sequence alignment and involves no heuristic.
The masker tokenises the \emph{canonical} SMILES and replaces whole BPE tokens
by single \texttt{<mask>} symbols, so re-tokenising the masked string cannot
change its length: an unmasked run is the same substring and re-merges
identically, and each \texttt{<mask>} is a hard boundary that prevents
neighbouring tokens merging across it. The masked string therefore has the same
token count as its canonical parent, and mask position $i$ holds what was parent
token $i$; this was verified at $100\%$ on 500 molecules, against only $16\%$
agreement when the raw rather than the canonical parent is used. The length
identity is asserted at every step and a pair that violates it is dropped rather
than trained on a silently misaligned target.

The fallback target is nonetheless a reconstruction target, and reconstructing
the parent attains $T_{\mathrm{c}}=1$, the worst possible value of the novelty
term of Eq.~\eqref{eq:stage10-loss}. Left at full weight it would make roughly
seven in ten updates train a copier. The scalar $\lambda_{\mathrm{fb}}$ (default
$0.3$) down-weights those updates alone, and spans the whole design space
between the two defensible extremes: $\lambda_{\mathrm{fb}}=1$ trains them as
ordinary examples, maximising the rate at which validity is acquired at the
greatest risk of collapse onto the parent, and $\lambda_{\mathrm{fb}}=0$ is
exactly equivalent to discarding those molecules.

\subsubsection{The differentiable objective}

With the target fixed, the gradient-carrying part of the stage is ordinary
masked-language-model cross-entropy on the logits of the original forward pass,
averaged over a molecule's masked positions so that a heavily masked molecule
does not dominate a mini-batch merely by having more of them:
\begin{equation}
\mathcal{L}_{\mathrm{CE}} \;=\; -\frac{1}{|\mathcal{M}|}
\sum_{i \,\in\, \mathcal{M}} \log P_{\theta}\bigl(t^{\star}_i \mid c(\tilde{s})\bigr).
\label{eq:stage10-ce}
\end{equation}
In variant~10a this is supplemented by an unlikelihood
term~\cite{welleck2020unlikelihood} that explicitly suppresses the tokens
observed to produce invalid molecules. Writing
\begin{align}
\mathcal{B} \;&=\; \bigl\{ (k,i) \;:\; \hat{s}^{(k)}\ \text{invalid},\;\;
i \in \mathcal{M},\;\; \hat{t}^{(k)}_i \neq t^{\star}_i \bigr\}, \notag \\
\mathcal{L}_{\mathrm{UL}} \;&=\; -\frac{1}{|\mathcal{B}|}
\sum_{(k,i) \,\in\, \mathcal{B}}
\log\Bigl(1 - P_{\theta}\bigl(\hat{t}^{(k)}_i \mid c(\tilde{s})\bigr)\Bigr),
\label{eq:stage10-ul}
\end{align}
the per-molecule objective is
\begin{equation}
\mathcal{L} \;=\; \omega\,\mathcal{L}_{\mathrm{CE}}
\;+\; \lambda_{\mathrm{UL}}\,\mathcal{L}_{\mathrm{UL}},
\label{eq:stage10-total}
\end{equation}
averaged over the molecules of a mini-batch. Variant~10b is
Eq.~\eqref{eq:stage10-total} with $\lambda_{\mathrm{UL}}=0$, that is,
cross-entropy towards the best candidate alone; the two variants were trained
from identical code and data into separate output directories, so that the
contribution of the unlikelihood term is isolated. Three details of
Eq.~\eqref{eq:stage10-ul} are deliberate. The penalty is written on
$1-P_{\theta}$ rather than as $+\log P_{\theta}$, because the latter diverges as
the probability is driven to zero, reintroducing exactly the unbounded direction
that motivated abandoning the policy gradient. Positions at which a rejected
candidate happens to agree with the target are excluded from $\mathcal{B}$,
since including them would set the two terms of Eq.~\eqref{eq:stage10-total}
against each other on the same token. And $\mathcal{L}_{\mathrm{UL}}$ is
\emph{not} scaled by $\omega$: on a molecule for which no candidate was valid,
sixteen decoded strings that failed to parse are legitimate evidence about what
should not be generated, unlike the parent reconstruction that $\omega$ exists
to discount. Gradients flow through $P_{\theta}$ alone; $\mathbf{t}^{\star}$,
$\mathcal{B}$ and every value of $\ell$ are constants of the step.

Informally, the chemistry acts as a teacher rather than as a signal. At each
step the model proposes sixteen ways of completing the same masked molecule;
RDKit reads all sixteen, discards those that are not molecules at all, ranks
what remains by drug-likeness, ease of synthesis, divergence from the parent and
absence of structural alerts, and hands back a single answer. The model is then
taught that answer in exactly the way it was pretrained --- ``at this masked
position, the token was this one'' --- and nothing about the chemistry needs to
be differentiable for that lesson to be delivered, because the only thing the
chemistry ever did was point. Validity improves not because a validity gradient
exists, but because the tokens repeatedly pointed at are the ones whose ring
closures balance, whose valences are legal and whose brackets match.

The same reading explains why this stage carries no Kullback--Leibler anchor to
the pretrained model, of the kind a policy-gradient formulation requires to stop
the policy drifting onto degenerate text. The targets of
Eq.~\eqref{eq:stage10-target} are either drawn from the model's own top-$k$
support or are the tokens of a real deposited molecule, so the objective cannot
move probability mass onto completions the model would never have proposed; the
anchoring is implicit in the construction of the target. The corresponding
weakness is the mirror image: an improvement the sampler never proposes is
unreachable, and with a $70\%$ fallback rate the majority of updates are parent
reconstructions at weight $\lambda_{\mathrm{fb}}$ rather than best-of-$K$
imitations. The trade-off that $\lambda_{\mathrm{fb}}$ governs --- validity
acquired by reconstruction against novelty lost to copying --- is therefore the
central hyperparameter of this stage.

\subsubsection{Learned toxicity term}

A fourth variant reinstates a learned toxicity estimate as a fifth term of the
selection loss,
\begin{equation}
\ell'(\hat{s}, s) \;=\; \ell(\hat{s}, s)
\;+\; w_{\mathrm{tox21}}\bigl(1 - \hat{p}_{\mathrm{clean}}(\hat{s})\bigr),
\qquad
\hat{p}_{\mathrm{clean}} \equiv 1 - \hat{p}_{\mathrm{tox21}},
\label{eq:stage104-loss}
\end{equation}
with the invalid branch likewise charged the full $w_{\mathrm{tox21}}$. The
worst attainable value for a valid candidate accordingly rises from $S$ to $S'
\equiv S + w_{\mathrm{tox21}}$ and an invalid candidate scores $S' +
w_{\mathrm{valid}}$, so the separating margin of
Eq.~\eqref{eq:stage10-ordering} remains exactly $w_{\mathrm{valid}}$ and
validity is still a lexicographic gate. The term was added on top of the
existing four weights rather than renormalising them, so that every other term
retains its meaning and the per-term diagnostics remain comparable across
variants; only the totals shift. $w_{\mathrm{tox21}} = 0.20$, matching
$w_{\mathrm{alert}}$, so that the rule-based and learned toxicity indicators
enter on equal footing and their disagreement on generated molecules is itself
observable. A candidate the classifier cannot score is charged the full weight,
the same pessimistic reading applied to every other term. One consequence is
worth stating because it would otherwise be silent: were no classifier
configured, the term would take the identical constant value for all $K$
candidates of a molecule and could not alter the arg-min of
Eq.~\eqref{eq:stage10-argmin}; the run would be numerically indistinguishable
from the variant without it while still reporting a toxicity curve, so the
implementation refuses to begin under that configuration rather than degrade
silently.

The classifier is a twelve-output sigmoid head fine-tuned from the same
ChemBERTa backbone as the generator on the Tox21 training
set~\cite{huang2016tox21,tox21challenge}, so that classifier and generator share
a tokenizer and the classifier sees the token stream the generator actually
emits. Of the $7{,}831 \times 12 = 93{,}972$ label cells, $77{,}946$ ($82.9\%$)
are observed; a blank denotes an assay not run rather than an inactive result,
so the training objective is a \emph{masked} binary cross-entropy under which
unobserved cells contribute neither loss nor gradient. Per-assay positive rates
span $2.9\%$ (NR-PPAR-$\gamma$) to $16.2\%$ (SR-ARE); each assay's positive
class is therefore reweighted by its training-split negative-to-positive ratio,
capped at $50$ so that no single rare assay dominates the gradient. Without that
reweighting the minimiser attains a near-optimal masked cross-entropy by
predicting inactivity everywhere, which would again render
$\hat{p}_{\mathrm{clean}}$ almost constant and the term inert. The classifier is
partitioned by the same framework rule applied to the generator, so that the two
splits mean the same thing, and its performance is reported as per-assay ROC-AUC
over the held-out frameworks.

\subsubsection{Optimisation and trainable parameters}

This stage fine-tunes ordinary weights rather than adapters: the LM head
together with the last $N=1$ of the six encoder blocks of
ChemBERTa-zinc-base-v1~\cite{chithrananda2020chemberta} was left trainable
($7.7$M of $44.1$M parameters, $17.5\%$) and the remainder frozen. Mini-batches
of 16 pairs were optimised with Adam~\cite{kingma2015adam} at learning rate
$5\times10^{-5}$ after clipping the global gradient norm to $1.0$, for 4 epochs
over the pooled, shuffled interaction-masked and randomly masked pairs. Only the
unfrozen tensors are written ($\approx\!30$\,MB against $180$\,MB for the whole
model), together with the optimiser moments and the state of every random
stream, so that an interrupted run resumes as the run it was rather than as a
new one begun from the same weights.

\subsubsection{Execution paths and the equivalence between them}

The objective above is specified once and executed by three interchangeable
implementations, so that the same weights can be produced either by the simplest
arrangement of the computation or by a throughput-oriented one that adapts to
whichever device a run lands on.

The reference arrangement performs one forward pass per molecule and one
cheminformatics measurement per candidate. The second dispatches the
$K\!\times\!B$ candidate measurements of a mini-batch to a persistent process
pool whose workers each pin their BLAS threads to one; because the measurement
is a pure function of two decoded strings, evaluating it in another process
cannot alter a floating-point result, and the pooled and serial objectives agree
exactly. This is the speedup that matters here: with $K=16$ candidates measured
per molecule, the cheminformatics term amounts to $256$ measurements per
optimiser update at the batch size used.

The third arrangement adds device-level optimisations: one right-padded $[B,L]$
forward pass per mini-batch, a single \texttt{multinomial} draw over the gathered
top-$k$ matrix for every masked position in the batch, bfloat16 or float16
autocast with the log-softmax and the cross-entropy retained in float32, TF32
for the remaining float32 matrix multiplications, length-bucketed batching,
thread and worker counts taken from the job's actual allocation rather than the
host core count, and DistributedDataParallel~\cite{li2020pytorchddp} with the
configured batch treated as a \emph{global} batch split across ranks, so that
the number of optimiser updates is independent of the number of devices.

These last optimisations are \emph{not} objective preserving, and the reason is
intrinsic to selection rather than incidental to the implementation. A score
that multiplies a log-probability tolerates a differently drawn sample: the
estimator remains unbiased. A score that \emph{selects} the target of
Eq.~\eqref{eq:stage10-argmin} does not. Three separate effects perturb the
candidate set --- drawing the batch in one \texttt{multinomial} call consumes the
random stream in a different order from one call per molecule; reduced precision
displaces the logits in their trailing bits; and dropout, sampled once per
forward pass, supplies $B$ independent masks to $B$ separate passes but only one
to a single padded pass. The third effect is the largest, displacing the logits
at masked positions by $\approx\!1.2$ on this model, which exceeds the
perturbation from reduced precision and cannot be removed by seeding, since the
two arrangements do not request the same shapes from the generator. Any of the
three can change which candidate attains the minimum in
Eq.~\eqref{eq:stage10-argmin}, and the model then learns a different token
sequence rather than a numerically perturbed update. The tuned arrangement is
therefore compared with the reference \emph{distributionally} --- over mean
candidate validity, mean composite loss and mean novelty --- and never by exact
output matching, while exact agreement is asserted in the deterministic limit of
$k=1$ with dropout disabled, which pins the batched indexing itself.

\subsubsection{Checkpointing and run provenance}

An epoch is not a cheap unit at $K=16$: on the capped training pool it is
several hours, so a checkpoint written only at epoch boundaries makes an
interruption expensive and a resumption inexact. The two resumable
implementations therefore write the unfrozen tensors, the Adam moments, the
running per-epoch aggregates, the batch cursor and the state of the Python, CPU
and CUDA random streams every fixed number of optimiser updates, by an atomic
replace, so that an interruption during the write cannot destroy the previous
state. Recording the optimiser is not a convenience: rebuilding Adam after each
restart discards the accumulated first and second moments and imposes a fresh
adaptation transient on every resumption, which would otherwise appear in the
training curve as an artefact of the interruption schedule rather than of the
objective. Recording the generator state is what makes a resumed run a
continuation of the run it interrupted, since the same stream draws the $K$
candidates of Eq.~\eqref{eq:stage10-sampling}.

Because the three implementations optimise one objective, a single run may
optionally be carried across them, beginning under the reference arrangement and
continuing under the tuned one as devices become available. Where that is done,
the checkpoint additionally carries a provenance record naming which
implementation trained which epochs, under which arithmetic precision and
sampling arrangement, and a resumption that changes any of these is reported at
the point it occurs. This is retained deliberately: by the argument above such a
run is not homogeneous in the way either implementation is on its own, the saved
weights carry no trace of the precision under which they were obtained, and the
distinction is not otherwise recoverable after the fact.

\subsubsection{Framework-based train--validation partition}

On the majority of molecules the training target is the parent's own tokens (the
fallback of Eq.~\eqref{eq:stage10-target}, taken on $\sim$70\% of molecules at
$p=15\%$). Reconstruction is therefore the modal training signal, which makes
memorisation of the training parents the specific failure mode this design
invites, and a property distribution computed over those same parents cannot
distinguish it from generalisation. The pooled training data were therefore
partitioned before any fine-tuning, and all held-out quantities are reported on
molecules the model has never been optimised towards.

The partition is made on \emph{molecular frameworks}, not on instances or on
molecules. Instance-level partitioning is untenable: a ligand deposited in many
structures would appear on both sides of the split tens of thousands of times.
Molecule-level partitioning removes that but leaves close analogues --- an
identical ring system bearing a different substituent --- free to straddle the
split, and such series are abundant in PDB-derived sets. Partitioning on the
Bemis--Murcko framework~\cite{bemis1996scaffolds}, the convention adopted by
MoleculeNet~\cite{wu2018moleculenet}, removes both: no validation molecule
shares a ring system with any training molecule. The consequence is that
held-out figures are systematically lower than the random-partition figures
conventionally quoted; that difference is the component of apparent performance
attributable to memorised chemistry, and its suppression is the purpose of the
exercise. Random partitions are known to overstate prospective performance for
exactly this reason~\cite{sheridan2013timesplit}, and ligand-based benchmarks
constructed without such control have been shown to reward memorisation over
generalisation~\cite{wallach2018memorization}.

Parents were canonicalised (RDKit~\cite{rdkit}) before grouping, so that a
molecule deposited under several equivalent notations forms one group rather
than several that could be assigned to opposite sides. Of the 46{,}982 distinct
parents in the interaction-masked source, 12{,}551 distinct frameworks were
obtained, and the three largest --- the acyclic class (6{,}754 molecules,
14.4\%), the Fe-porphyrin/haem framework (4{,}089, 8.7\%) and the
Mg-porphyrin/chlorophyll framework (2{,}393, 5.1\%) --- together hold 28\% of
the data.

How such groups are dealt out therefore matters as much as how they are formed.
Groups are shuffled uniformly and assigned to validation until the fold's
molecule budget is met, \emph{taking} the group that crosses the budget rather
than skipping it. Skipping is what would make size matter: a group could only be
held out if it fitted, so the largest families could never be, and the
validation fold would consist of a handful of scaffold families whose effective
sample size is the number of families rather than the number of molecules.
Taking the crossing group costs an overshoot of at most one group and buys the
property the estimate rests on --- a group's chance of being held out does not
depend on how big it is. A residual dependence remains and is second-order: a
group still perturbs the budget it is measured against, so its inclusion
probability differs from a singleton's by roughly its own share of all
molecules, about 3\% relative for the largest group here. Size-blindness does
not imply low variance: drawing one of the three large families into validation
makes it alone 15--60\% of the fold, which is a legitimate draw, so the split
report states the largest held-out group's share explicitly and warns beyond
25\%. A further 1{,}015 parents (2.2\%), whose SMILES exceed 400 characters and
are predominantly peptide and oligosaccharide ligands, bypass the framework
reduction and form singleton groups keyed on canonical SMILES; these are
partitioned at molecule granularity, the conservative fallback.

The partition is computed once over the complete parent set and cached as a
manifest read by every variant. It is therefore independent of the per-parent
instance cap, of any sub-sampling of the training pool and of the variant being
run, so a molecule cannot change fold between runs and all implementations agree
on what is held out. A molecule absent from the manifest inherits the fold of
its framework, and one whose framework is also absent defaults to training --- a
direction that can never inflate a held-out score. Validation molecules were
evaluated at the end of each epoch under the identical best-of-$K$ procedure
with gradients disabled; the epoch attaining the lowest mean validation loss was
retained as the final model; and property distributions are reported separately
for training and held-out molecules, the difference between them being the
memorisation measurement that within-training diagnostics cannot supply.

\subsubsection{Diagnostics}

Four quantities were logged per epoch and are the intended read-out of whether
the objective is behaving as constructed: the mean of
Eq.~\eqref{eq:stage10-total}; the mean of $\ell$ at the selected target; the
\emph{candidate validity rate}, the fraction of all $K$ sampled completions that
parse; and the \emph{fallback rate}, the fraction of molecules with
$\mathcal{V}=0$. The last two are the diagnostic pair of interest. A run in
which candidate validity rises while the novelty term of
Eq.~\eqref{eq:stage10-loss} rises alongside it is one that is purchasing
validity by memorising parents, and calls for a smaller $\lambda_{\mathrm{fb}}$;
the intended behaviour is a falling fallback rate with the novelty term flat or
falling.

\begin{table}[htbp]
\centering
\caption{Default hyperparameters for property-guided supervised fine-tuning.}
\label{tab:stage10-hparams}
\begin{tabular}{@{}ll@{}}
\toprule
Parameter & Value \\
\midrule
Base model & ChemBERTa-zinc-base-v1~\cite{chithrananda2020chemberta} \\
Trainable parameters & LM head $+$ last $N=1$ of 6 encoder blocks (7.7M / 44.1M, 17.5\%) \\
Mask rate $p$ (\% of BPE tokens) & 15\% \\
Decoding & one-shot (all masks in a single forward pass) \\
Candidates per molecule $K$ & 16 \\
Sampling top-$k$ / temperature $T$ & 20 / 1.2 \\
Loss weights $(w_{\mathrm{QED}}, w_{\mathrm{SA}}, w_{\mathrm{nov}}, w_{\mathrm{alert}})$ & (0.30, 0.20, 0.30, 0.20), summing to $S=1$ \\
Invalidity margin $w_{\mathrm{valid}}$ & 1.0 (an invalid candidate scores $S+w_{\mathrm{valid}}=2$) \\
Parent-fallback weight $\lambda_{\mathrm{fb}}$ & 0.3 \\
Unlikelihood weight $\lambda_{\mathrm{UL}}$ & 0.5 (variant 10a); 0 (variant 10b) \\
Learned toxicity weight $w_{\mathrm{tox21}}$ & not used by default; $0.20$ in the toxicity-aware variant \\
Worst valid / invalid loss & $S=1$, $S+w_{\mathrm{valid}}=2$; with toxicity: $S'=1.2$, $2.2$ \\
Train--validation partition & Bemis--Murcko framework, size-blind, 10\% held out, seed 42 \\
Model selection & epoch of lowest mean validation loss \\
Batch size & 16 \\
Optimiser / learning rate & Adam / $5\times10^{-5}$ \\
Gradient-norm clip & 1.0 \\
Epochs & 4 \\
\bottomrule
\end{tabular}
\end{table}

\FloatBarrier

\section{Results}
To evaluate the efficacy of the interaction-guided masking strategy, we developed a preprocessing pipeline that integrates structural protein–ligand data from the Protein–Ligand Interaction Profiler (PLIP) with the ChemBERTa transformer architecture. By selectively masking "attractive" atoms—those not participating in critical non-covalent interactions such as hydrogen bonds and hydrophobic contacts—we demonstrated that the model can be guided to reconstruct molecular sequences that prioritize biologically relevant regions.


In figure 1  Pink-shaded atoms labelled with an asterisk~($*$) are the
PLIP-identified interaction sites selected for ChemBERTa token
masking; interaction types follow the legend inset in each panel
(piStacking: dark green dash-dot; hydrogen bond: dark blue solid;
water bridge: cyan solid; hydrophobic contact: light grey dashed).
\textbf{(a)}~Ligand 1GH (PDB: 3ZYU) engages the BRD4 acetyl-lysine
pocket through the richest polar network of the five systems: one
$\pi$-stacking contact with residue~81A from the bicyclic
imidazolone ring, two water bridges from the isoxazole N$^{*}$ and
imidazolone N$^{*}$, one hydrogen bond from the pyridine N, and five
hydrophobic contacts distributed across the bicyclic
oxazole--purine scaffold.
The dense pink masking pattern reflects the high density of
interaction-critical atoms on this compact scaffold.
\textbf{(b)}~Ligand 62G (PDB: 5HLS) presents a simpler polar
profile: one hydrogen bond from the pendant amide N--H to a backbone
carbonyl, one water bridge from the isoxazole N$^{*}$, and seven
hydrophobic contacts spanning the benzodiazepine--isoxazole core; no
$\pi$-stacking is detected.
Masked atoms are concentrated on the bicyclic aromatic portion and
the amide carbonyl carbon, highlighting the pharmacophoric relevance
of these positions.
\textbf{(c)}~Ligand EAM (PDB: 3P5O) shows one hydrogen bond to the
backbone amide N$^{*}$ and one water bridge from the same region,
together with six hydrophobic contacts including a chlorine-mediated
contact at the para-chlorophenyl moiety of the tricyclic
benzodiazepine core; no $\pi$-stacking is detected.
The Cl atom participates directly in hydrophobic burial, providing a
halogen-mediated pharmacophoric anchor captured by the masking scheme.
\textbf{(d)}~Ligand JQ1 (PDB: 4QZS) engages BRD4 through two water
bridges anchored at the triazole N$^{*}$ atoms and eight hydrophobic
contacts distributed across the thienodiazepine--triazole scaffold,
including contacts at the thiophene S$^{*}$ and the methyl-decorated
diazepine C$^{*}$ positions; no hydrogen bond or $\pi$-stacking is
detected in this crystal environment.
\textbf{(e)}~Ligand JQ1 (PDB: 3MXF) presents a distinct interaction
profile: the two water bridges observed in 4QZS are absent and
replaced by a single hydrogen bond to the triazole N, alongside eight
hydrophobic contacts on the same thienodiazepine--triazole scaffold.
The shift from water-mediated to direct hydrogen bonding between the
two JQ1 co-crystals reflects subtle differences in binding pose and
active-site water occupancy, yielding distinct PLIP-ranked masking
patterns for each crystal form.
In all panels, the masked atoms~($*$) define the interaction-aware
(NIP Masking ) token positions fed to ChemBERTa for conditional SMILES
generation in subsequent pipeline stages.%
\normalsize

%                      -
%                      -
%  FIGURE 1  —  panels + one-line title caption
%                      -
\begin{figure}[t!]
    \centering
 
    %  - Row 1: (a) 1GH  and  (b) 62G  -
    \begin{subfigure}[b]{0.38\textwidth}
        \centering
        \includegraphics[width=\textwidth]{{vis_inta/1GH_A_1173_masked.2d_interactions.png}}
        \caption{\scriptsize 2D Interaction: ligand 1GH (PDB: 3ZYU)}
        \label{fig:audit_1gh}
    \end{subfigure}
    \hfill
    \begin{subfigure}[b]{0.38\textwidth}
        \centering
        \includegraphics[width=\textwidth]%
            {vis_inta/62G_A_201_masked.2d_interactions.png}
        \caption{\scriptsize 2D Interaction: ligand 62G (PDB: 5HLS)}
        \label{fig:audit_62g}
    \end{subfigure}
 
    \vspace{4pt}
 
    %  - Row 2: (c) EAM  and  (d) JQ1 (4QZS)  -
    \begin{subfigure}[b]{0.38\textwidth}
        \centering
        \includegraphics[width=\textwidth]%
            {vis_inta/EAM_A_1_masked.2d_interactions.png}
        \caption{\scriptsize 2D Interaction: ligand EAM (PDB: 3P5O)}
        \label{fig:audit_eam}
    \end{subfigure}
    \hfill
    \begin{subfigure}[b]{0.38\textwidth}
        \centering
        \includegraphics[width=\textwidth]%
            {vis_inta/JQ1_A_1_masked.2d_interactions.png}
        \caption{\scriptsize 2D Interaction: ligand JQ1 (PDB: 4QZS)}
        \label{fig:audit_jq1_4qzs}
    \end{subfigure}
 
    \vspace{4pt}
 
    %  - Row 3: (e) JQ1 (3MXF) centred  -
    \begin{subfigure}[b]{0.38\textwidth}
        \centering
        \includegraphics[width=\textwidth]%
            {vis_inta/JQ1_A_201_masked.2d_interactions.png}
        \caption{\scriptsize 2D Interaction: ligand JQ1 (PDB: 3MXF)}
        \label{fig:audit_jq1_3mxf}
    \end{subfigure}
 
    %  - Title-only caption (full legend lives below)  -
    \caption{%
        \textbf{PLIP-based 2D interactions of five BRD4--ligand
        complexes in ATTRACTIVE mode.}
    }
    \label{fig:plip_2d_audit}
\end{figure}
 
%                      -
%  FIGURE LEGEND  (Nature style: placed in text flow immediately
%  after \end{figure}, not inside the float)
%                      -

 
\
 


%                    -
%  FIGURE — Incremental Generation Curves (5 panels, 2+2+1)
%                    -
\begin{figure}[htbp]
    \centering
 
    %  - Row 1: (a) 1GH  and  (b) 62G  -
    \begin{subfigure}[b]{0.47\textwidth}
        \centering
        \includegraphics[width=\textwidth]{incremental_curves/incremental curves non interactive/1GH-A-1173.png}
        \caption{\scriptsize Ligand 1GH (PDB: 3ZYU)}
        \label{fig:igc_1gh}
    \end{subfigure}
    \hfill
    \begin{subfigure}[b]{0.47\textwidth}
        \centering
        \includegraphics[width=\textwidth]{incremental_curves/incremental curves non interactive/62G-A-201.png}
        \caption{\scriptsize Ligand 62G (PDB: 5HLS)}
        \label{fig:igc_62g}
    \end{subfigure}
 
    \vspace{0.5em}
 
    %  - Row 2: (c) JQ1-A-1  and  (d) JQ1-A-201  -
    \begin{subfigure}[b]{0.47\textwidth}
        \centering
        \includegraphics[width=\textwidth]{incremental_curves/incremental curves non interactive/JQ1-A-1.png}
        \caption{\scriptsize Ligand JQ1 (PDB: 4QZS)}
        \label{fig:igc_jq1_4qzs}
    \end{subfigure}
    \hfill
    \begin{subfigure}[b]{0.47\textwidth}
        \centering
        \includegraphics[width=\textwidth]{incremental_curves/incremental curves non interactive/JQ1-A-201.png}
        \caption{\scriptsize Ligand JQ1 (PDB: 3MXF)}
        \label{fig:igc_jq1_3mxf}
    \end{subfigure}
 
    \vspace{0.5em}
 
    %  - Row 3: (e) EAM centred  -
    \begin{subfigure}[b]{0.47\textwidth}
        \centering
        \includegraphics[width=\textwidth]{incremental_curves/incremental curves non interactive/EAM-A-1.png}
        \caption{\scriptsize Ligand EAM (PDB: 3P5O)}
        \label{fig:igc_eam}
    \end{subfigure}
 
    %  - Full-width caption  -
    \caption{%
        \textbf{Incremental generation curves for five BRD4--ligand complexes. Comparison of Random vs Non-interacting Part Masking }
        Each panel reports the number of unique, chemically valid SMILES
        produced by ChemBERTa at each masking depth (mask\_count $= 1,\ldots,N$),
        Comparing two atom-selection strategies:
        \textit{Interaction-Aware} (NIP Masking , blue solid line), in which atoms are
        chosen for masking in order of their PLIP-detected interaction rank,
        and \textit{Random masking} (red dashed line), in which atoms are
        selected uniformly at random.
        %
        (\textbf{a})~Ligand 1GH exhibits a sharp NIP Masking peak at mask\_count~$= 2$
        ($\approx$48 molecules) with a secondary lobe at mask\_count~$= 6$
        ($\approx$24), while random masking collapses below 10 molecules by
        mask\_count~$= 3$, demonstrating that interaction-guided site selection
        substantially amplifies generative yield for this compact scaffold.
        %
        (\textbf{b})~Ligand 62G shows an anomalous inversion: random masking
        reaches a pronounced peak at mask\_count~$= 3$ ($\approx$159),
        exceeding the NIP Masking peak ($\approx$74); this may reflect the high
        symmetry and topological redundancy of the benzodiazepine--isoxazole
        scaffold, which makes any single-atom perturbation productive
        regardless of interaction rank.
        %
        (\textbf{c,d})~Both JQ1 co-crystal structures (PDB: 4QZS and 3MXF)
        display higher random peaks at mask\_count~$\leq 3$ but sustained NIP Masking 
        generation at larger mask counts, indicating that interaction-aware
        masking covers a wider productive chemical neighbourhood around the
        thienodiazepine--triazole pharmacophore.
        %
        (\textbf{e})~Ligand EAM shows the most consistent NIP Masking advantage:
        NIP Masking outperforms random at every tested depth, peaks at mask\_count~$= 3$
        ($\approx$39), and maintains non-zero yield through mask\_count~$= 7$,
        whereas random generation decays monotonically after mask\_count~$= 2$.
        %
        In all panels, generation yield approaches zero beyond a
        compound-specific saturation threshold, beyond which the masked
        SMILES fragment becomes too large for the pretrained ChemBERTa
        vocabulary to reconstruct a valid molecule.
    }
    \label{fig:incremental_generation_curves}
\end{figure}





  




 
%                      
%  Column-header helpers (no extra package needed)
%                      
\newcommand{\colhdr}[1]{%
    {\centering\small\textbf{#1}\par}\vspace{2pt}%
}
 

% ================================================================
%  PAGE 1  —  Rows 1–3  (panels a–f)
% ================================================================
\begin{figure}[p]
    \centering
 
   %  - Column headers  -
    \begin{minipage}[t]{0.44\textwidth}
        \centering\small\textbf{Non-interacting part Masking}
    \end{minipage}
    \hfill
    \begin{minipage}[t]{0.44\textwidth}
        \centering\small\textbf{Random Masking}
    \end{minipage}
     
    \vspace{2pt}
 
    % ==== ROW 1 : 1GH-A-1173 ====
    \begin{subfigure}[b]{0.44\textwidth}
        \centering
        \includegraphics[width=\textwidth]{scatter_plots_CNN_composite/non_interation_part_masked/scatter_1GH-A-1173_composite.png}
        \caption{\scriptsize 1GH-A-1173 — NIP Masking ($n = 101$)}
        \label{fig:scatter_1gh_ia}
    \end{subfigure}
    \hfill
    \begin{subfigure}[b]{0.44\textwidth}
        \centering
        \includegraphics[width=\textwidth]{scatter_plots_CNN_composite/scatter_1GH-A-1173_rand.png}
        \caption{\scriptsize 1GH-A-1173 — Random ($n = 21$)}
        \label{fig:scatter_1gh_rand}
    \end{subfigure}
 
    \vspace{2pt}
 
    % ==== ROW 2 : EAM-A-1 ====
    \begin{subfigure}[b]{0.44\textwidth}
        \centering
        \includegraphics[width=\textwidth]{scatter_plots_CNN_composite/non_interation_part_masked/scatter_EAM-A-1_composite.png}
        \caption{\scriptsize EAM-A-1 — NIP Masking ($n = 126$)}
        \label{fig:scatter_eam_ia}
    \end{subfigure}
    \hfill
    \begin{subfigure}[b]{0.44\textwidth}
        \centering
        \includegraphics[width=\textwidth]{scatter_plots_CNN_composite/scatter_EAM-A-1_rand.png}
        \caption{\scriptsize EAM-A-1 — Random ($n = 92$)}
        \label{fig:scatter_eam_rand}
    \end{subfigure}
 
    \vspace{2pt}
 
    % ==== ROW 3 : JQ1-A-201 ====
    \begin{subfigure}[b]{0.44\textwidth}
        \centering
        \includegraphics[width=\textwidth]{scatter_plots_CNN_composite/non_interation_part_masked/scatter_JQ1-A-201_composite.png}
        \caption{\scriptsize JQ1-A-201 — NIP Masking ($n = 210$)}
        \label{fig:scatter_jq1201_ia}
    \end{subfigure}
    \hfill
    \begin{subfigure}[b]{0.44\textwidth}
        \centering
        \includegraphics[width=\textwidth]{scatter_plots_CNN_composite/scatter_JQ1-A-201_rand.png}
        \caption{\scriptsize JQ1-A-201 — Random ($n = 223$)}
        \label{fig:scatter_jq1201_rand}
    \end{subfigure}
 
    %  - Page-1 caption  -
    \caption{%
        \textbf{Novelty--potency trade-off for interaction-aware (NIP Masking ) versus
        random masking across five BRD4 ligand systems (page~1 of~2).}
        Each scatter plot reports all successfully docked molecules generated
        by ChemBERTa under the respective masking strategy.
        The \textit{x}-axis encodes Tanimoto similarity to the cognate
        reference ligand; molecules to the left of the dashed line
        (Tanimoto~$= 0.40$) are classified as \emph{novel-scaffold}
        candidates by the standard medicinal-chemistry threshold.
        The \textit{y}-axis reports the composite docking score
        $S = \text{CNNaffinity} \times \text{CNNscore}$, where higher values
        indicate stronger predicted binding~\cite{ragoza2017protein}.
        Rows are ordered from the largest NIP Masking yield advantage (top) to
        near-parity (bottom); see page~2 for conditions where random masking
        dominates and the cross-ligand summary.
        %
        (\textbf{a,b})~For 1GH-A-1173, NIP Masking generates \textasciitilde5$\times$
        more molecules than random (101 vs.\ 21), with substantial coverage
        of the novel-scaffold region (Tanimoto~$= 0.18$--$0.39$) and
        competitive scores up to $S \approx 4.2$ in that zone; random
        masking produces almost exclusively analogue-like molecules
        (Tanimoto~$\geq 0.60$), demonstrating that interaction-guided site
        selection is essential for exploratory generation on this compact
        bicyclic scaffold.
        %
        (\textbf{c,d})~EAM-A-1 under NIP Masking yields 37\% more molecules (126
        vs.\ 92) and achieves the highest individual composite score across
        all conditions ($S \approx 7.3$); however, all NIP Masking -generated
        molecules reside at Tanimoto~$\geq 0.38$, whereas random masking
        reaches down to Tanimoto~$\approx 0.20$, indicating that PLIP-ranked
        interaction sites on the rigid tricyclic benzodiazepine core bias
        generation toward structurally similar derivatives.
        %
        (\textbf{e,f})~JQ1-A-201 shows near-parity between strategies
        (210 vs.\ 223 molecules); NIP Masking populates the Tanimoto~$= 0.35$--$0.45$
        transition zone more densely, suggesting that interaction-aware
        masking can identify lead-hopping candidates at the scaffold
        boundary.
    }
    \label{fig:NIP Masking _vs_rand_docking}
\end{figure}
 
 
% ================================================================
%  PAGE 2  —  Rows 4–6  (panels g–l)
%  \ContinuedFloat keeps the same figure number and continues
%  the subfigure alphabet (g, h, i, j, k, l)
% ================================================================
\begin{figure}[p]
    \ContinuedFloat          % <-- continues Fig. \ref{fig:NIP Masking _vs_rand_docking}
    \centering
 
    %  - Column headers (repeat for reader orientation)  -
   %  - Column headers  -
    \begin{minipage}[t]{0.44\textwidth}
        \centering\small\textbf{Non-interacting part Masking}
    \end{minipage}
    \hfill
    \begin{minipage}[t]{0.44\textwidth}
        \centering\small\textbf{Random Masking}
    \end{minipage}
 
    \vspace{2pt}
 
    % ==== ROW 4 : JQ1-A-1 ====
    \begin{subfigure}[b]{0.44\textwidth}
        \centering
        \includegraphics[width=\textwidth]{scatter_plots_CNN_composite/non_interation_part_masked/scatter_JQ1-A-1_composite.png}
        \caption{\scriptsize JQ1-A-1 — NIP Masking ($n = 123$)}
        \label{fig:scatter_jq1_ia}
    \end{subfigure}
    \hfill
    \begin{subfigure}[b]{0.44\textwidth}
        \centering
        \includegraphics[width=\textwidth]{scatter_plots_CNN_composite/scatter_JQ1-A-1_rand.png}
        \caption{\scriptsize JQ1-A-1 — Random ($n = 183$)}
        \label{fig:scatter_jq1_rand}
    \end{subfigure}
 
    \vspace{2pt}
 
    % ==== ROW 5 : 62G-A-201 ====
    \begin{subfigure}[b]{0.44\textwidth}
        \centering
        \includegraphics[width=\textwidth]{scatter_plots_CNN_composite/non_interation_part_masked/scatter_62G-A-201_composite.png}
        \caption{\scriptsize 62G-A-201 — NIP Masking ($n = 172$)}
        \label{fig:scatter_62g_ia}
    \end{subfigure}
    \hfill
    \begin{subfigure}[b]{0.44\textwidth}
        \centering
        \includegraphics[width=\textwidth]{scatter_plots_CNN_composite/scatter_62G-A-201_rand.png}
        \caption{\scriptsize 62G-A-201 — Random ($n = 283$)}
        \label{fig:scatter_62g_rand}
    \end{subfigure}
 
    \vspace{2pt}
 
    % ==== ROW 6 : All ligands combined (summary row) ====
    \begin{subfigure}[b]{0.44\textwidth}
        \centering
        \includegraphics[width=\textwidth]{scatter_plots_CNN_composite/non_interation_part_masked/scatter_combined_composite.png}
        \caption{\scriptsize All ligands combined — NIP Masking ($n = 732$)}
        \label{fig:scatter_combined_ia}
    \end{subfigure}
    \hfill
    \begin{subfigure}[b]{0.44\textwidth}
        \centering
        \includegraphics[width=\textwidth]{scatter_plots_CNN_composite/scatter_combined_rand.png}
        \caption{\scriptsize All ligands combined — Random ($n = 802$)}
        \label{fig:scatter_combined_rand}
    \end{subfigure}
 
    %  - Page-2 caption  -
    \caption{%
        \textbf{(Continued from previous page.)}
        %
        (\textbf{g,h})~JQ1-A-1 is the first condition where random masking
        yields substantially more molecules (183 vs.\ 123, a 49\%
        surplus), with a dense high-scoring cluster at
        Tanimoto~$= 0.80$--$1.0$ ($S > 7.0$); NIP Masking achieves similar maximum
        scores but lower total yield, reflecting that the
        thienodiazepine--triazole pharmacophore contains chemically tolerant
        positions not fully captured by the top-ranked PLIP sites alone.
        %
        (\textbf{i,j})~62G-A-201 shows the strongest random yield advantage
        (283 vs.\ 172, 65\% more molecules); however, a disproportionate
        fraction of random-generated compounds fall in the low-score band
        ($S = 1.0$--$2.5$), whereas NIP Masking -generated molecules concentrate
        between $S = 3.0$ and $6.5$, indicating a quality-over-quantity
        advantage for NIP Masking on the benzodiazepine--isoxazole scaffold.
        %
        (\textbf{k,l})~The cross-ligand summary panels (NIP Masking : $n = 732$;
        random: $n = 802$) confirm that NIP Masking produces a more structured
        distribution with clear inter-ligand clustering; EAM-A-1 (green)
        contributes the highest-scoring candidates under both strategies.
        The random combined panel shows a broader, more diffuse cloud
        shifted toward higher Tanimoto values, consistent with a portfolio
        that is chemically less diverse overall.
    }
    % No new \label here; the figure number is shared with page 1
\end{figure}
% ================================================================
%  Vinardo Novelty–Potency Scatter Plots
%  PAGE 1  —  Rows 1–3  (panels a–f)
% ================================================================
\begin{figure}[p]
    \centering

    %  - Column headers  -
    %  - Column headers  -
    \begin{minipage}[t]{0.44\textwidth}
        \centering\small\textbf{Non-interacting part Masking}
    \end{minipage}
    \hfill
    \begin{minipage}[t]{0.44\textwidth}
        \centering\small\textbf{Random Masking}
    \end{minipage}

    \vspace{2pt}

    % ==== ROW 1 : 1GH-A-1173 ====
    \begin{subfigure}[b]{0.44\textwidth}
        \centering
        \includegraphics[width=\textwidth]{vinardo_docking/scatter_1GH-A-1173.png}
        \caption{\scriptsize 1GH-A-1173   NIP Masking ($n = 98$)}
        \label{fig:vin_scatter_1gh_ia}
    \end{subfigure}
    \hfill
    \begin{subfigure}[b]{0.44\textwidth}
        \centering
        \includegraphics[width=\textwidth]{vinardo_docking/scatter_1GH-A-1173_rand.png}
        \caption{\scriptsize 1GH-A-1173   Random ($n = 21$)}
        \label{fig:vin_scatter_1gh_rand}
    \end{subfigure}

    \vspace{2pt}

    % ==== ROW 2 : EAM-A-1 ====
    \begin{subfigure}[b]{0.44\textwidth}
        \centering
        \includegraphics[width=\textwidth]{vinardo_docking/scatter_EAM-A-1.png}
        \caption{\scriptsize EAM-A-1   NIP Masking ($n = 148$)}
        \label{fig:vin_scatter_eam_ia}
    \end{subfigure}
    \hfill
    \begin{subfigure}[b]{0.44\textwidth}
        \centering
        \includegraphics[width=\textwidth]{vinardo_docking/scatter_EAM-A-1_rand.png}
        \caption{\scriptsize EAM-A-1   Random ($n = 92$)}
        \label{fig:vin_scatter_eam_rand}
    \end{subfigure}

    \vspace{2pt}

    % ==== ROW 3 : JQ1-A-201 ====
    \begin{subfigure}[b]{0.44\textwidth}
        \centering
        \includegraphics[width=\textwidth]{vinardo_docking/scatter_JQ1-A-201.png}
        \caption{\scriptsize JQ1-A-201   NIP Masking ($n = 177$)}
        \label{fig:vin_scatter_jq1201_ia}
    \end{subfigure}
    \hfill
    \begin{subfigure}[b]{0.44\textwidth}
        \centering
        \includegraphics[width=\textwidth]{vinardo_docking/scatter_JQ1-A-201_rand.png}
        \caption{\scriptsize JQ1-A-201   Random ($n = 223$)}
        \label{fig:vin_scatter_jq1201_rand}
    \end{subfigure}

    \caption{%
        \textbf{Novelty--potency trade-off for interaction-aware (NIP Masking ) versus
        random masking across five BRD4 ligand systems, Vinardo scoring
        (page~1 of~2).}
        Each scatter plot reports all successfully docked molecules generated
        by ChemBERTa under the respective masking strategy.
        The \textit{x}-axis encodes Tanimoto similarity to the cognate
        reference ligand (Morgan fingerprints, radius~2); molecules to the
        left of the dashed line (Tanimoto~$= 0.40$) are classified as
        \emph{novel-scaffold} candidates by the standard medicinal-chemistry
        threshold.
        The \textit{y}-axis reports the Vinardo empirical docking score
        (kcal/mol); lower (more negative) values indicate stronger
        predicted binding affinity~\cite{quiroga2016vinardo}.
        Rows are ordered from the largest NIP Masking yield advantage (top) to
        near-parity (bottom); see page~2 for conditions where random
        masking dominates and the cross-ligand summary.
        %
        (\textbf{a,b})~For 1GH-A-1173, NIP Masking generates \textasciitilde5$\times$
        more molecules than random (98 vs.\ 21), with coverage of the
        novel-scaffold region (Tanimoto~$< 0.40$) and competitive Vinardo
        scores reaching $\approx -9.3$\,kcal/mol; random masking produces
        almost exclusively analogue-like molecules (Tanimoto~$\geq 0.60$),
        confirming that interaction-guided site selection is essential for
        exploratory generation on this compact bicyclic scaffold.
        %
        (\textbf{c,d})~EAM-A-1 under NIP Masking yields 61\% more molecules (148
        vs.\ 92) and achieves the strongest Vinardo scores
        ($\approx -10.0$\,kcal/mol); NIP masking molecules span a wide Tanimoto
        range including novel-scaffold candidates, whereas random masking
        concentrates output at high similarity with weaker scores overall.
        %
        (\textbf{e,f})~JQ1-A-201 shows near-parity in yield (177 vs.\ 223);
        NIP Masking populates the transition zone (Tanimoto~$= 0.35$--$0.45$) more
        densely, suggesting interaction-aware masking can identify
        lead-hopping candidates at the scaffold boundary.
    }
    \label{fig:NIP Masking _vs_rand_vinardo}
\end{figure}


% ================================================================
%  Vinardo Novelty–Potency Scatter Plots
%  PAGE 2  —  Rows 4–6  (panels g–l)
% ================================================================
\begin{figure}[p]
    \ContinuedFloat
    \centering

    %  - Column headers (repeat for reader orientation)  -
   \begin{minipage}{0.44\textwidth}
        \centering\small\textbf{Non Interacting Part Masking}
    \end{minipage}
    \hfill
    \begin{minipage}{0.44\textwidth}
        \centering\small\textbf{Random Masking}
    \end{minipage}
    \vspace{2pt}

    % ==== ROW 4 : JQ1-A-1 ====
    \begin{subfigure}[b]{0.44\textwidth}
        \centering
        \includegraphics[width=\textwidth]{vinardo_docking/scatter_JQ1-A-1.png}
        \caption{\scriptsize JQ1-A-1   NIP Masking ($n = 196$)}
        \label{fig:vin_scatter_jq1_ia}
    \end{subfigure}
    \hfill
    \begin{subfigure}[b]{0.44\textwidth}
        \centering
        \includegraphics[width=\textwidth]{vinardo_docking/scatter_JQ1-A-1_rand.png}
        \caption{\scriptsize JQ1-A-1   Random ($n = 183$)}
        \label{fig:vin_scatter_jq1_rand}
    \end{subfigure}

    \vspace{2pt}

    % ==== ROW 5 : 62G-A-201 ====
    \begin{subfigure}[b]{0.44\textwidth}
        \centering
        \includegraphics[width=\textwidth]{vinardo_docking/scatter_62G-A-201.png}
        \caption{\scriptsize 62G-A-201   NIP Masking ($n = 231$)}
        \label{fig:vin_scatter_62g_ia}
    \end{subfigure}
    \hfill
    \begin{subfigure}[b]{0.44\textwidth}
        \centering
        \includegraphics[width=\textwidth]{vinardo_docking/scatter_62G-A-201_rand.png}
        \caption{\scriptsize 62G-A-201   Random ($n = 283$)}
        \label{fig:vin_scatter_62g_rand}
    \end{subfigure}

    \vspace{2pt}

    % ==== ROW 6 : All ligands combined (summary row) ====
    \begin{subfigure}[b]{0.44\textwidth}
        \centering
        \includegraphics[width=\textwidth]{vinardo_docking/scatter_combined.png}
        \caption{\scriptsize All ligands combined   NIP Masking ($n = 850$)}
        \label{fig:vin_scatter_combined_ia}
    \end{subfigure}
    \hfill
    \begin{subfigure}[b]{0.44\textwidth}
        \centering
        \includegraphics[width=\textwidth]{vinardo_docking/scatter_combined_rand.png}
        \caption{\scriptsize All ligands combined   Random ($n = 802$)}
        \label{fig:vin_scatter_combined_rand}
    \end{subfigure}

    \caption{%
        \textbf{(Continued from previous page.)}
        %
        (\textbf{g,h})~JQ1-A-1 is the first condition where NIP Masking yields
        more molecules than random (196 vs.\ 183); both strategies produce
        strong Vinardo scores in the $-8$ to $-9.5$\,kcal/mol range,
        with NIP Masking achieving slightly more coverage in the novel-scaffold
        region (Tanimoto~$< 0.40$).
        %
        (\textbf{i,j})~62G-A-201 shows the strongest random yield
        advantage (283 vs.\ 231); however, random-generated molecules
        include a disproportionate fraction of low-scoring compounds,
        whereas NIP Masking -generated molecules concentrate in the stronger
        binding range, indicating a quality-over-quantity advantage
        for NIP Masking on this scaffold.
        %
        (\textbf{k,l})~The cross-ligand summary panels (NIP Masking : $n = 850$;
        random: $n = 802$) confirm that NIP Masking produces a more structured
        distribution with clear inter-ligand clustering at competitive
        Vinardo scores; the random combined panel shows a broader,
        more diffuse cloud shifted toward higher Tanimoto values and
        weaker predicted binding affinities overall.
    }
    % No new \label here; the figure number is shared with page 1
\end{figure}
%========== Tanimato Samilarity =======
 
% ================================================================
%  PAGE 1  —  Rows 1–3  (panels a–f)
% ================================================================
\begin{figure}[p]
    \centering
 
    %  - Column headers  -
    %  - Column headers  -
    \begin{minipage}[t]{0.44\textwidth}
        \centering\small\textbf{Non-interacting part Masking}
    \end{minipage}
    \hfill
    \begin{minipage}[t]{0.44\textwidth}
        \centering\small\textbf{Random Masking}
    \end{minipage}
 
    \vspace{2pt}
 
    % ==== ROW 1 : 1GH-A-1173  (NIP Masking mean 0.372 — most diverse) ====
    \begin{subfigure}[b]{0.44\textwidth}
        \centering
        \includegraphics[width=\textwidth]{Tanimoto_sim/1GH-A-1173/histogram_pairwise_ia.png}
        \caption{\scriptsize 1GH-A-1173 — NIP Masking 
                 ($\bar{T}=0.372$, $n=122$, $N=12$)}
        \label{fig:hist_1gh_ia}
    \end{subfigure}
    \hfill
    \begin{subfigure}[b]{0.44\textwidth}
        \centering
        \includegraphics[width=\textwidth]{Tanimoto_sim/1GH-A-1173/histogram_pairwise_rand.png}
        \caption{\scriptsize 1GH-A-1173 — Random
                 ($\bar{T}=0.539$, $n=21$, $N=12$)}
        \label{fig:hist_1gh_rand}
    \end{subfigure}
 
    \vspace{2pt}
 
    % ==== ROW 2 : JQ1-A-201  (NIP masking mean 0.420) ====
    \begin{subfigure}[b]{0.44\textwidth}
        \centering
        \includegraphics[width=\textwidth]{Tanimoto_sim/JQ1-A-201/histogram_pairwise_ia.png}
        \caption{\scriptsize JQ1-A-201 — NIP Masking 
                 ($\bar{T}=0.420$, $n=210$, $N=10$)}
        \label{fig:hist_jq1201_ia}
    \end{subfigure}
    \hfill
    \begin{subfigure}[b]{0.44\textwidth}
        \centering
        \includegraphics[width=\textwidth]{Tanimoto_sim/JQ1-A-201/histogram_pairwise_rand.png}
        \caption{\scriptsize JQ1-A-201 — Random
                 ($\bar{T}=0.528$, $n=229$, $N=10$)}
        \label{fig:hist_jq1201_rand}
    \end{subfigure}
 
    \vspace{2pt}
 
    % ==== ROW 3 : 62G-A-201  (NIP masking mean 0.477) ====
    \begin{subfigure}[b]{0.44\textwidth}
        \centering
        \includegraphics[width=\textwidth]{Tanimoto_sim/62G-A-201/histogram_pairwise_ia.png}
        \caption{\scriptsize 62G-A-201 — NIP Masking 
                 ($\bar{T}=0.477$, $n=172$, $N=10$)}
        \label{fig:hist_62g_ia}
    \end{subfigure}
    \hfill
    \begin{subfigure}[b]{0.44\textwidth}
        \centering
        \includegraphics[width=\textwidth]{Tanimoto_sim/62G-A-201/histogram_pairwise_rand.png}
        \caption{\scriptsize 62G-A-201 — Random
                 ($\bar{T}=0.417$, $n=283$, $N=10$)}
        \label{fig:hist_62g_rand}
    \end{subfigure}
 
    %  - Page-1 caption  -
    \caption{%
        \textbf{Pairwise Tanimoto similarity distributions of
        ChemBERTa-generated molecules under interaction-aware (NIP Masking )
        and random masking across five BRD4 ligand systems
        (page~1 of~2).}
        %
        Each histogram aggregates all pairwise Tanimoto similarity
        scores computed over the full set of unique valid molecules
        generated across all masking depths ($N$ masks) for the
        respective ligand and strategy.
        The red dashed line marks the mean pairwise similarity
        $\bar{T}$; lower $\bar{T}$ indicates higher internal
        chemical diversity within the generated library.
        Rows are ordered by ascending NIP masking mean Tanimoto (most
        internally diverse library first).
        %
        (\textbf{a,b})~For 1GH-A-1173 (PDB:~3ZYU), NIP Masking produces a
        library of $n = 122$ molecules with a notably low mean
        pairwise similarity ($\bar{T}_{\mathrm{NIP Masking }} = 0.372$,
        right-skewed distribution), indicating high internal
        structural diversity; the random library ($n = 21$,
        $\bar{T}_{\mathrm{rand}} = 0.539$) is far smaller and
        considerably less diverse, with most pairs clustering
        above Tanimoto~$= 0.60$.
        %
        (\textbf{c,d})~JQ1-A-201 (PDB:~3MXF) shows a moderate NIP Masking
        diversity ($\bar{T}_{\mathrm{NIP Masking}} = 0.420$, $n = 210$)
        with a smooth unimodal distribution centred near
        Tanimoto~$= 0.40$; the random library
        ($\bar{T}_{\mathrm{rand}} = 0.528$, $n = 229$) has a
        broader, flatter distribution extending into the high
        similarity tail (Tanimoto~$\geq 0.70$), suggesting the
        random strategy recovers a greater proportion of
        near-identical analogues.
        %
        (\textbf{e,f})~62G-A-201 (PDB:~5HLS) is the only system
        where random masking yields \emph{lower} mean pairwise
        similarity than NIP Masking
        ($\bar{T}_{\mathrm{rand}} = 0.417$ vs.\
        $\bar{T}_{\mathrm{NIP Masking}} = 0.477$), driven by a large random
        library ($n = 283$) with a sharp modal peak at
        Tanimoto~$\approx 0.42$ and a heavy low-similarity tail;
        the NIP Masking distribution is sharper and more concentrated,
        consistent with PLIP-guided masking targeting a
        constrained set of interaction-critical positions on the
        benzodiazepine--isoxazole scaffold.
    }
    \label{fig:pairwise_tanimoto}
\end{figure}
 
 
% ================================================================
%  PAGE 2  —  Rows 4–5  (panels g–j)
%  \ContinuedFloat preserves the figure number and continues
%  the subfigure alphabet from (g)
% ================================================================
\begin{figure}[p]
    \ContinuedFloat
    \centering
 
    %  - Column headers (repeated for reader orientation)  -
       %  - Column headers  -
    \begin{minipage}[t]{0.44\textwidth}
        \centering\small\textbf{Non-interacting part Masking}
    \end{minipage}
    \hfill
    \begin{minipage}[t]{0.44\textwidth}
        \centering\small\textbf{Random Masking}
    \end{minipage}
 
    \vspace{2pt}
 
    % ==== ROW 4 : JQ1-A-1  (NIP masking mean 0.486) ====
    \begin{subfigure}[b]{0.44\textwidth}
        \centering
        \includegraphics[width=\textwidth]{Tanimoto_sim/JQ1-A-1/histogram_pairwise_ia.png}
        \caption{\scriptsize JQ1-A-1 — NIP masking
                 ($\bar{T}=0.486$, $n=123$, $N=11$)}
        \label{fig:hist_jq1_ia}
    \end{subfigure}
    \hfill
    \begin{subfigure}[b]{0.44\textwidth}
        \centering
        \includegraphics[width=\textwidth]{Tanimoto_sim/JQ1-A-1/histogram_pairwise_rand.png}
        \caption{\scriptsize JQ1-A-1 — Random
                 ($\bar{T}=0.504$, $n=185$, $N=11$)}
        \label{fig:hist_jq1_rand}
    \end{subfigure}
 
    \vspace{2pt}
 
    % ==== ROW 5 : EAM-A-1  (NIP masking mean 0.534 — least diverse) ====
    \begin{subfigure}[b]{0.44\textwidth}
        \centering
        \includegraphics[width=\textwidth]{Tanimoto_sim/EAM-A-1/histogram_pairwise_ia.png}
        \caption{\scriptsize EAM-A-1 — NIP masking
                 ($\bar{T}=0.534$, $n=126$, $N=8$)}
        \label{fig:hist_eam_ia}
    \end{subfigure}
    \hfill
    \begin{subfigure}[b]{0.44\textwidth}
        \centering
        \includegraphics[width=\textwidth]{Tanimoto_sim/EAM-A-1/histogram_pairwise_rand.png}
        \caption{\scriptsize EAM-A-1 — Random
                 ($\bar{T}=0.475$, $n=92$, $N=8$)}
        \label{fig:hist_eam_rand}
    \end{subfigure}
 
    %  - Page-2 caption  -
    \caption{%
        \textbf{(Continued from previous page.)}
        %
        (\textbf{g,h})~JQ1-A-1 (PDB:~4QZS) exhibits nearly
        identical mean pairwise similarities between strategies
        ($\bar{T}_{\mathrm{NIP masking }} = 0.486$ vs.\
        $\bar{T}_{\mathrm{rand}} = 0.504$, $\Delta\bar{T} = 0.018$),
        the smallest ia--random difference across all five systems;
        however, the random library is 50\% larger ($n = 185$
        vs.\ $n = 123$) and shows a prominent high-similarity
        secondary peak at Tanimoto~$\approx 0.60$, indicating that
        random masking of the thienodiazepine--triazole scaffold
        produces a bimodal distribution mixing novel variants and
        near-identical analogues, whereas NIP masking  yields a
        narrower, unimodal profile.
        %
        (\textbf{i,j})~EAM-A-1 (PDB:~3P5O) is the only system
        where NIP masking produces a \emph{higher} mean pairwise
        similarity than random ($\bar{T}_{\mathrm{NIP Masking}} = 0.534$
        vs.\ $\bar{T}_{\mathrm{rand}} = 0.475$), yet generates
        more molecules ($n = 126$ vs.\ $n = 92$); the NIP Masking
        distribution is tightly bell-shaped centred near
        Tanimoto~$= 0.50$, whereas the random distribution is
        bimodal with a low-similarity peak near
        Tanimoto~$= 0.25$ and a high-similarity peak near
        Tanimoto~$= 0.75$, suggesting that random masking on the
        rigid benzodiazepine core samples two disconnected regions
        of chemical space rather than a coherent neighbourhood.
        %
        Taken together, the five-system comparison shows that NIP
        masking consistently produces more unimodal, coherent
        similarity distributions, whereas random masking more
        frequently yields bimodal or heavy-tailed profiles that
        include a higher proportion of near-identical molecules.
    }
    % No \label on the continuation — label shared with page 1
\end{figure}
 

 
% ================================================================
%  FIGURE A — Nearest-Neighbour Similarity Histograms
%  PAGE 1  :  rows 1–3  (panels a–f)
% ================================================================
\begin{figure}[p]
    \centering
 
      %  - Column headers  -
    \begin{minipage}[t]{0.44\textwidth}
        \centering\small\textbf{Non-interacting part Masking}
    \end{minipage}
    \hfill
    \begin{minipage}[t]{0.44\textwidth}
        \centering\small\textbf{Random Masking}
    \end{minipage}
 
    \vspace{2pt}
 
    % ==== ROW 1 : 1GH-A-1173 ====
    \begin{subfigure}[b]{0.44\textwidth}
        \centering
        \includegraphics[width=\textwidth]%
            {chembel_ana/similarity_score_histogram_1GH-A-1173_ia.png}
        \caption{\scriptsize 1GH-A-1173 — NIP Masking}
        \label{fig:nn_sim_1gh_ia}
    \end{subfigure}
    \hfill
    \begin{subfigure}[b]{0.44\textwidth}
        \centering
        \includegraphics[width=\textwidth]%
            {chembel_ana/similarity_score_histogram_1GH-A-1173_rand.png}
        \caption{\scriptsize 1GH-A-1173 — Random}
        \label{fig:nn_sim_1gh_rand}
    \end{subfigure}
 
    \vspace{2pt}
 
    % ==== ROW 2 : JQ1-A-201 ====
    \begin{subfigure}[b]{0.44\textwidth}
        \centering
        \includegraphics[width=\textwidth]%
            {chembel_ana/similarity_score_histogram_JQ1-A-201_ia.png}
        \caption{\scriptsize JQ1-A-201 — NIP Masking}
        \label{fig:nn_sim_jq1201_ia}
    \end{subfigure}
    \hfill
    \begin{subfigure}[b]{0.44\textwidth}
        \centering
        \includegraphics[width=\textwidth]%
            {chembel_ana/similarity_score_histogram_JQ1-A-201_rand.png}
        \caption{\scriptsize JQ1-A-201 — Random}
        \label{fig:nn_sim_jq1201_rand}
    \end{subfigure}
 
    \vspace{2pt}
 
    % ==== ROW 3 : 62G-A-201 ====
    \begin{subfigure}[b]{0.44\textwidth}
        \centering
        \includegraphics[width=\textwidth]%
            {chembel_ana/similarity_score_histogram_62G-A-201_ia.png}
        \caption{\scriptsize 62G-A-201 — NIP Masking}
        \label{fig:nn_sim_62g_ia}
    \end{subfigure}
    \hfill
    \begin{subfigure}[b]{0.44\textwidth}
        \centering
        \includegraphics[width=\textwidth]%
            {chembel_ana/similarity_score_histogram_62G-A-201_rand.png}
        \caption{\scriptsize 62G-A-201 — Random}
        \label{fig:nn_sim_62g_rand}
    \end{subfigure}
 
    \caption{%
        \textbf{Nearest-neighbour Tanimoto similarity of generated
        molecules to the ChEMBL database, interaction-aware (NIP Masking)
        vs.\ random masking — page~1 of~2.}
        Each histogram shows, for every generated molecule, the
        Tanimoto similarity to its single closest neighbour in the
        ChEMBL compound collection (Morgan fingerprints, radius~2).
        The red dashed line at $T = 0.40$ marks the standard
        scaffold-similarity threshold; molecules to the right are
        considered analogues of known ChEMBL compounds, while those
        to the left represent genuinely novel chemical matter.
        %
        (\textbf{a,b})~For 1GH-A-1173, NIP Masking generates a bimodal
        distribution with a novel-scaffold peak near $T \approx 0.30$
        and a broader analogue population at $T = 0.45$--$0.65$,
        demonstrating simultaneous exploration of both novel and
        known chemical space; the random library (only $n = 21$) is
        sparse and concentrated at $T \geq 0.75$, indicating that
        random masking almost exclusively recovers near-identical
        ChEMBL compounds.
        %
        (\textbf{c,d})~JQ1-A-201 NIP Masking produces a unimodal distribution
        centred near $T \approx 0.45$, with most molecules classified
        as scaffold analogues; random masking yields a prominent
        bimodal distribution with a second peak at $T \approx 0.80$,
        reflecting a greater proportion of high-similarity ChEMBL
        matches in the random library.
        %
        (\textbf{e,f})~For 62G-A-201, both strategies concentrate
        above the $T = 0.40$ threshold; the NIP Masking peak at
        $T \approx 0.55$ is sharper and lower in median similarity
        than the random peak at $T \approx 0.45$--$0.50$, which is
        broader and extends further into the high-similarity tail.
        See page~2 for JQ1-A-1 and EAM-A-1.
    }
    \label{fig:nn_similarity_hist}
\end{figure}
 
 
% ================================================================
%  FIGURE A — PAGE 2  :  rows 4–5  (panels g–j)
% ================================================================
\begin{figure}[p]
    \ContinuedFloat
    \centering
 
    %  - Column headers  -
       %  - Column headers  -
    \begin{minipage}[t]{0.44\textwidth}
        \centering\small\textbf{Non-interacting part Masking}
    \end{minipage}
    \hfill
    \begin{minipage}[t]{0.44\textwidth}
        \centering\small\textbf{Random Masking}
    \end{minipage}
 
    % ==== ROW 4 : JQ1-A-1 ====
    \begin{subfigure}[b]{0.44\textwidth}
        \centering
        \includegraphics[width=\textwidth]%
            {chembel_ana/similarity_score_histogram_JQ1-A-1_ia.png}
        \caption{\scriptsize JQ1-A-1 — NIP Masking}
        \label{fig:nn_sim_jq1_ia}
    \end{subfigure}
    \hfill
    \begin{subfigure}[b]{0.44\textwidth}
        \centering
        \includegraphics[width=\textwidth]%
            {chembel_ana/similarity_score_histogram_JQ1-A-1_rand.png}
        \caption{\scriptsize JQ1-A-1 — Random}
        \label{fig:nn_sim_jq1_rand}
    \end{subfigure}
 
    \vspace{2pt}
 
    % ==== ROW 5 : EAM-A-1 ====
    \begin{subfigure}[b]{0.44\textwidth}
        \centering
        \includegraphics[width=\textwidth]%
            {chembel_ana/similarity_score_histogram_EAM-A-1_ia.png}
        \caption{\scriptsize EAM-A-1 — NIP Masking}
        \label{fig:nn_sim_eam_ia}
    \end{subfigure}
    \hfill
    \begin{subfigure}[b]{0.44\textwidth}
        \centering
        \includegraphics[width=\textwidth]%
            {chembel_ana/similarity_score_histogram_EAM-A-1_rand.png}
        \caption{\scriptsize EAM-A-1 — Random}
        \label{fig:nn_sim_eam_rand}
    \end{subfigure}
 
    \caption{%
        \textbf{(Continued from previous page.)}
        %
        (\textbf{g,h})~JQ1-A-1 shows a clear divergence between
        strategies: NIP Masking generates a unimodal distribution centred at
        $T \approx 0.60$ with a well-defined analogue population
        in the medium-similarity range; random masking produces a
        strongly right-skewed distribution with a sharp peak at
        $T \approx 0.80$, indicating a systematic bias toward
        high-similarity ChEMBL matches and reduced chemical novelty
        relative to NIP Masking.
        %
        (\textbf{i,j})~EAM-A-1 presents the most striking contrast:
        NIP Masking yields a right-skewed distribution with a mode at
        $T \approx 0.75$--$0.80$, consistent with the tight
        interaction-constrained neighbourhood of the tricyclic
        benzodiazepine core; random masking generates a bimodal
        distribution with a novel-scaffold peak near
        $T \approx 0.30$ and a high-similarity peak near
        $T \approx 0.85$, suggesting that random perturbations
        of the rigid scaffold sample two disconnected regions of
        ChEMBL chemical space.
        Across all five systems, NIP Masking consistently narrows
        the nearest-neighbour similarity distribution toward
        a coherent mid-range neighbourhood ($T = 0.40$--$0.70$),
        whereas random masking more frequently populates the
        high-similarity tail ($T \geq 0.80$), reflecting a
        tendency to recover near-identical ChEMBL compounds rather
        than structurally diverse analogues.
    }
    % No \label — shared with page 1
\end{figure}
 
 
% ================================================================
%  FIGURE B — Nearest-Neighbour Frequency Bar Charts
%  PAGE 1  :  rows 1–3  (panels a–f)
% ================================================================
\begin{figure}[p]
    \centering
 
    %  - Column headers  -
   \begin{minipage}{0.44\textwidth}
        \centering\small\textbf{Non interacting part Masking}
    \end{minipage}
    \hfill
    \begin{minipage}{0.44\textwidth}
        \centering\small\textbf{Random Masking}
    \end{minipage}
 
    \vspace{2pt}
 
    % ==== ROW 1 : 1GH-A-1173 ====
    \begin{subfigure}[b]{0.44\textwidth}
        \centering
        \includegraphics[width=\textwidth]%
            {chembel_ana/nearest_neighbour_frequency_1GH-A-1173_ia.png}
        \caption{\scriptsize 1GH-A-1173 — NIP Masking}
        \label{fig:nn_freq_1gh_ia}
    \end{subfigure}
    \hfill
    \begin{subfigure}[b]{0.44\textwidth}
        \centering
        \includegraphics[width=\textwidth]%
            {chembel_ana/nearest_neighbour_frequency_1GH-A-1173_rand.png}
        \caption{\scriptsize 1GH-A-1173 — Random}
        \label{fig:nn_freq_1gh_rand}
    \end{subfigure}
 
    \vspace{2pt}
 
    % ==== ROW 2 : JQ1-A-201 ====
    \begin{subfigure}[b]{0.44\textwidth}
        \centering
        \includegraphics[width=\textwidth]%
            {chembel_ana/nearest_neighbour_frequency_JQ1-A-201_ia.png}
        \caption{\scriptsize JQ1-A-201 — NIP Masking}
        \label{fig:nn_freq_jq1201_ia}
    \end{subfigure}
    \hfill
    \begin{subfigure}[b]{0.44\textwidth}
        \centering
        \includegraphics[width=\textwidth]%
            {chembel_ana/nearest_neighbour_frequency_JQ1-A-201_rand.png}
        \caption{\scriptsize JQ1-A-201 — Random}
        \label{fig:nn_freq_jq1201_rand}
    \end{subfigure}
 
    \vspace{2pt}
 
    % ==== ROW 3 : 62G-A-201 ====
    \begin{subfigure}[b]{0.44\textwidth}
        \centering
        \includegraphics[width=\textwidth]%
            {chembel_ana/nearest_neighbour_frequency_62G-A-201_ia.png}
        \caption{\scriptsize 62G-A-201 — NIP Masking}
        \label{fig:nn_freq_62g_ia}
    \end{subfigure}
    \hfill
    \begin{subfigure}[b]{0.44\textwidth}
        \centering
        \includegraphics[width=\textwidth]%
            {chembel_ana/nearest_neighbour_frequency_62G-A-201_rand.png}
        \caption{\scriptsize 62G-A-201 — Random}
        \label{fig:nn_freq_62g_rand}
    \end{subfigure}
 
    \caption{%
        \textbf{Nearest-neighbour frequency of generated molecules
        mapped to ChEMBL compounds ($T \geq 0.40$), interaction-aware
        (NIP Masking) vs.\ random masking — page~1 of~2.}
        Each bar represents a ChEMBL compound that is the nearest
        neighbour ($T \geq 0.40$) of at least one generated molecule;
        bar height gives the count of generated molecules mapping to
        that compound; stacked colours denote similarity range
        (yellow: Low, $0.40$--$0.60$; light green: Medium,
        $0.60$--$0.80$; dark green: High, $0.80$--$1.00$);
        annotated median similarity (Med) is shown above each bar.
        Total count sorts bars in descending order.
        %
        (\textbf{a,b})~For 1GH-A-1173, NIP masking maps 65 generated
        molecules to CHEMBL2181820 (Med:~0.62) with a diverse
        similarity range spanning Low to High; random masking maps
        only 11 molecules to the same top hit, but at a higher median
        similarity (Med:~0.85), confirming that NIP Masking explores a broader
        and a more chemically diverse neighbourhood of this reference
        compound.
        %
        (\textbf{c,d})~JQ1-A-201 NIP masking maps 146 generated molecules to Non-interacting part Masking
        CHEMBL1957266 (Med:~0.55), predominantly in the Low
        similarity range; random masking maps 174 molecules to the
        same compound but at a substantially higher median
        (Med:~0.78), With the High similarity tier dominating,
        indicating that random masking produces analogues that are
        structurally closer to this known JQ1-like scaffold.
        %
        (\textbf{e,f})~62G-A-201 NIP Masking identifies CHEMBL4303404 (PDB:
        5HLS, the reference compound itself, Med:~0.56) as the
        dominant nearest neighbour ($n \approx 148$); random masking
        maps more molecules to this compound ($n \approx 154$) but
        also recovers a large secondary cluster at CHEMBL5274793
        ($n \approx 75$, Med:~0.64), reflecting the broader chemical
        space sampled by random masking on this scaffold.
        See page~2 for JQ1-A-1 and EAM-A-1.
    }
    \label{fig:nn_frequency}
\end{figure}
 
 
% ================================================================
%  FIGURE B — PAGE 2:  rows 4–5  (panels g–j)
% ================================================================
\begin{figure}[p]
    \ContinuedFloat
    \centering
 
    %  - Column headers  -
       %  - Column headers  -
    \begin{minipage}[t]{0.44\textwidth}
        \centering\small\textbf{Non-interacting part Masking}
    \end{minipage}
    \hfill
    \begin{minipage}[t]{0.44\textwidth}
        \centering\small\textbf{Random Masking}
    \end{minipage}
 
    \vspace{2pt}
 
    % ==== ROW 4 : JQ1-A-1 ====
    \begin{subfigure}[b]{0.44\textwidth}
        \centering
        \includegraphics[width=\textwidth]%
            {chembel_ana/nearest_neighbour_frequency_JQ1-A-1_ia.png}
        \caption{\scriptsize JQ1-A-1 — NIP Masking}
        \label{fig:nn_freq_jq1_ia}
    \end{subfigure}
    \hfill
    \begin{subfigure}[b]{0.44\textwidth}
        \centering
        \includegraphics[width=\textwidth]%
            {chembel_ana/nearest_neighbour_frequency_JQ1-A-1_rand.png}
        \caption{\scriptsize JQ1-A-1 — Random}
        \label{fig:nn_freq_jq1_rand}
    \end{subfigure}
 
    \vspace{2pt}
 
    % ==== ROW 5 : EAM-A-1 ====
    \begin{subfigure}[b]{0.44\textwidth}
        \centering
        \includegraphics[width=\textwidth]%
            {chembel_ana/nearest_neighbour_frequency_EAM-A-1_ia.png}
        \caption{\scriptsize EAM-A-1 — NIP Masking}
        \label{fig:nn_freq_eam_ia}
    \end{subfigure}
    \hfill
    \begin{subfigure}[b]{0.44\textwidth}
        \centering
        \includegraphics[width=\textwidth]%
            {chembel_ana/nearest_neighbour_frequency_EAM-A-1_rand.png}
        \caption{\scriptsize EAM-A-1 — Random}
        \label{fig:nn_freq_eam_rand}
    \end{subfigure}
 
    \caption{%
        \textbf{(Continued from previous page.)}
        %
        (\textbf{g,h})~JQ1-A-1 NIP masking maps the largest absolute number
        of molecules to CHEMBL1957266 across all five systems
        ($n \approx 230$, Med:~0.65), with a broad similarity
        distribution spanning all three tiers; random masking maps
        135 molecules to the same compound at a higher median
        (Med:~0.70), and recovers fewer secondary neighbours,
        indicating that NIP Masking generates a richer chemical neighbourhood
        around this central BRD4-active scaffold.
        %
        (\textbf{i,j})~EAM-A-1 NIP masking maps 56 generated molecules to
        CHEMBL1232461 (PDB: 7AQT, Med:~0.74) with a medium- to
        high-similarity profile; random masking maps 47 molecules
        to the same compound at a higher median (Med:~0.85),
        with the High tier dominant.
        The NIP Masking library recovers a wider panel of secondary ChEMBL
        neighbours (8 compounds vs.\ 5 for random), consistent
        with broader pharmacophoric exploration enabled by
        PLIP-guided masking of the tricyclic benzodiazepine core.
        %
        Across both panels and all five systems, NIPmasking
        consistently distributes generated molecules across a
        larger number of distinct ChEMBL nearest neighbours with
        lower mediaMn similarity scores, indicating wider chemical
        space coverage relative to random masking, which
        concentrates output around fewer, higher-similarity
        reference compounds.
    }
    % No \label — shared with page 1
\end{figure}


% ================================================================
%  FIGURE 7 — ChEMBL Nearest-Neighbour Tanimoto Boxplot Summary
%  NIP Masking vs. Random masking — all five ligand groups
% ================================================================
\begin{figure}[htbp]
    \centering

  \begin{minipage}{0.44\textwidth}
        \centering\small\textbf{Non interacting part Masking}
    \end{minipage}
    \hfill
    \begin{minipage}{0.44\textwidth}
        \centering\small\textbf{Random Masking}
    \end{minipage}

    \vspace{2pt}

    \begin{subfigure}[b]{0.44\textwidth}
        \centering
        \includegraphics[width=\textwidth]{summary plot_stage 5/boxplot_chembl_tanimoto_ia.png}
        \caption{\scriptsize All ligands — NIP masking}
        \label{fig:boxplot_chembl_ia}
    \end{subfigure}
    \hfill
    \begin{subfigure}[b]{0.44\textwidth}
        \centering
        \includegraphics[width=\textwidth]{summary plot_stage 5/boxplot_chembl_tanimoto_rand.png}
        \caption{\scriptsize All ligands — Random masking}
        \label{fig:boxplot_chembl_rand}
    \end{subfigure}

    \caption{%
        \textbf{ChEMBL nearest-neighbour Tanimoto similarity across all
        five BRD4 ligand groups,  non interaction vs.\ random
        masking   Stage~5 summary ($N = 5$ groups, all nearest-neighbour
        scores).}
        Each box represents the distribution of per-molecule nearest-neighbour
        Tanimoto scores to the ChEMBL compound collection (Morgan fingerprints,
        radius~2) for all valid generated molecules within each ligand group
        (red line: mediaMn; box: interquartile range; whiskers:
        $1.5\times\mathrm{IQR}$; circles: outliers).
        The grey dashed reference line at $T = 0.40$ marks the conventional
        scaffold-similarity threshold; all mediaMn values across both strategies
        lie above this threshold, confirming that the generated libraries are
        pharmacologically anchored to known bioactive chemical space.
        %
        (\textbf{a})~Under ia masking, 1GH-A-1173 returns the highest median
        nearest-neighbour similarity ($T \approx 0.76$) with a narrow
        interquartile range, reflecting the tight PLIP-constrained neighbourhood
        of its compact bicyclic scaffold; the four remaining groups cluster
        between $T = 0.54$ and $T = 0.64$, with JQ1-A-1 showing the widest
        IQR, consistent with the larger and more chemically diverse NIP Masking library
        generated for that scaffold.
        %
        (\textbf{b})~Under random masking, the five groups converge toward a
        more uniform distribution centred in the $T = 0.55$--$0.67$ range,
        with broader interquartile ranges for 1GH-A-1173 and EAM-A-1
        indicating greater between-molecule variability in ChEMBL coverage;
        notably, EAM-A-1 displays a cluster of low-similarity outliers near
        $T \approx 0.15$--$0.21$, reflecting the bimodal ChEMBL coverage
        seen in the per-molecule histograms (Fig.~\ref{fig:nn_similarity_hist}).
        Across both panels, NIP masking yields tighter, more consistent
        nearest-neighbour distributions than random masking, confirming that
        PLIP-guided site selection produces a pharmacologically coherent chemical
        libraries across structurally diverse BRD4 scaffolds.
    }
    \label{fig:boxplot_chembl_summary}
\end{figure}

% ================================================================
%  METHODS — ALGORITHM 1: Tanimoto-Cutoff Sequential Filter
% ================================================================
\clearpage
\paragraph{Tanimoto-cutoff sequent NIP Masking filter:}
Generated molecules from the non-interaction-site masking pool
were ranked by a four-step sequential NIP Masking filter applied in strict
order (Algorithm~\ref{alg:idea1}).
The filter first removes physically implausible binding poses
using two GNINA-derived quality gates: a CNN-based pose-plausibility
threshold (CNNscore~$> 0.80$) and a physics-based affinity
threshold (Vinardo score~$< -8.0$~kcal/mol~\cite{quiroga2016vinardo}).
Only molecules satisfying both gates are considered further.
A hard scaffold-novelty criterion (Tanimoto similarity to the
reference ligand~$< 0.40$; Morgan fingerprints, radius~$r = 2$,
$L = 2048$~bits) then removes analogues of the reference
compound~\cite{maggiora2014}.
The surviving molecules are sorted by composite score
$S = \text{CNNaffinity} \times \text{CNNscore}$ in descending
order and the top-5 are reported (Supplementary Table~S1
provides the full top-10).
 
\begin{algorithm}[htbp]
\caption{Tanimoto-Cutoff Sequential Filter}
\label{alg:idea1}
\begin{algorithmic}[1]
\Require Docked molecule set $\mathcal{G}$; thresholds
         $\tau_{\text{CNN}} = 0.80$,
         $\tau_{\text{Vin}} = -8.0$~kcal/mol,
         $\delta = 0.40$
\Ensure  Top-5 novel potent candidates $\mathcal{C}_1$
\State $\mathcal{G}_1 \gets
       \{m \in \mathcal{G} \mid \text{CNNscore}(m) > \tau_{\text{CNN}}\}$
       \Comment{Remove poorly-posed compounds}
\State $\mathcal{G}_2 \gets
       \{m \in \mathcal{G}_1 \mid \text{Vinardo}(m) < \tau_{\text{Vin}}\}$
       \Comment{Remove weak physical binders}
\State $\mathcal{G}_3 \gets
       \{m \in \mathcal{G}_2 \mid
       T(\mathbf{f}_m,\mathbf{f}_{\text{ref}}) < \delta\}$
       \Comment{Enforce scaffold novelty}
\State $\mathcal{G}_3^{\star} \gets
       \text{sort}\,(\mathcal{G}_3,\;
       S_{\text{composite}}\downarrow)$
       \Comment{Rank by $S = \text{CNNaffinity} \times \text{CNNscore}$}
\State \Return $\text{top-5}(\mathcal{G}_3^{\star})$
\end{algorithmic}
\end{algorithm}
 
 
% ================================================================
%  METHODS — ALGORITHM 2: Butina Scaffold Clustering
% ================================================================
 
\paragraph{Butina scaffold-cluster representative
selection:}
To ensure that selected candidates span the broadest possible
range of distinct chemical scaffolds present in the
non-interaction-masked library, a scaffold-cluster
representative strategy was implemented using the Butina
single-linkage clustering algorithm~\cite{butina1999}
(Algorithm~\ref{alg:idea5}).
A symmetric pairwise Tanimoto distance matrix
$d_{ij} = 1 - T(\mathbf{f}_i, \mathbf{f}_j)$ was computed
for all $n$ generated molecules using Morgan fingerprints
(radius~$r = 2$, $L = 2048$~bits).
The Butina algorithm partitioned the library into clusters
such that all members of a cluster share a Tanimoto similarity
$\geq 1 - \varepsilon$ to the cluster centroid
(default $\varepsilon = 0.40$, consistent with the Bemis--Murcko
scaffold-equivalence threshold~\cite{bemis1996scaffolds}).
From each cluster the molecule with the highest composite score
was selected as the sole representative, and all representatives
were sorted by composite score to yield the final ranked list.
The top-5 are reported; Supplementary Table~S2 provides the
full top-10.
 
\begin{algorithm}[htbp]
\caption{Butina Scaffold-Cluster Representative
         Selection}
\label{alg:idea5}
\begin{algorithmic}[1]
\Require Docked molecule set $\mathcal{G}$; cluster cutoff
         $\varepsilon = 0.40$
\Ensure  Top-5 scaffold-diverse potent candidates $\mathcal{C}_5$
\For{each pair $(i,j)$, $i > j$}
    \State $d_{ij} \gets 1 - T(\mathbf{f}_i, \mathbf{f}_j)$
           \Comment{Tanimoto distance matrix (upper triangle)}
\EndFor
\State $\{K_1, K_2, \ldots, K_C\} \gets
       \text{Butina}(\{d_{ij}\},\, \varepsilon)$
       \Comment{Partition into $C$ scaffold clusters}
\For{each cluster $K_c$, $c = 1,\ldots,C$}
    \State $m_c^{*} \gets
           \operatorname{arg\,max}_{m \in K_c}
           S_{\text{composite}}(m)$
           \Comment{Best-composite representative}
\EndFor
\State $\mathcal{R} \gets
       \{m_c^{*}\}_{c=1}^{C}$
       sorted by $S_{\text{composite}}$ descending
\State \Return $\text{top-5}(\mathcal{R})$
\end{algorithmic}
\end{algorithm}
 


\begin{table*}[ht]
\centering
\caption{Tanimoto-cutoff sequential filtering}
\label{tab:generated_molecules}
\renewcommand{\arraystretch}{1.3}
\begin{adjustbox}{width=\textwidth}
\begin{tabular}{|
>{\centering\arraybackslash}m{2.6cm}|
>{\raggedright\arraybackslash}p{6.2cm}|
c|c|c|c|c|}
\hline
\textbf{Structure} &
\textbf{SMILES} &
\shortstack{\textbf{CNNaffinity}\\(higher is better)} &
\shortstack{\textbf{CNNscore}\\(higher is better)} &
\shortstack{\textbf{Vinardo}\\(lower is better)} &
\shortstack{\textbf{Tanimoto}\\(lower = more novel)} &
\shortstack{\textbf{Composite}\\(aff $\times$ score)} \\
\hline
%       -- Molecule 1       --
\begin{adjustbox}{max width=2.4cm, max height=2.4cm, center}
\includegraphics{composite_cnn/1.png}
\end{adjustbox}
&
{\tiny\seqsplit{[Cl]c1[cH][cH][c](C2=CCCCC[C@@H]([CH2][C](=O)[NH][CH2][CH3])CSc32[n]cc([CH3])c3-[c]3[c]2[cH]c([O]2)[cH]c3)[cH][cH]1}}
&
7.4031 & 0.9144 & -9.0633 & 0.3131 & 6.7692 \\
\hline
%       -- Molecule 2       --
\begin{adjustbox}{max width=2.4cm, max height=2.4cm, center}
\includegraphics{composite_cnn/2.png}
\end{adjustbox}
&
{\tiny\seqsplit{[n]1c([CH3])CCCc2[c]([o]1)[C@H]([CH2][C]([NH2])=[O])[N]=[C]([c]1[cH][cH]c([Cl])[cH][cH]1)c1[cH][cH][cH]CS1-2}}
&
6.9855 & 0.9520 & -9.1952 & 0.3974 & 6.6503 \\
\hline
%       -- Molecule 3       --
\begin{adjustbox}{max width=2.4cm, max height=2.4cm, center}
\includegraphics{composite_cnn/3.png}
\end{adjustbox}
&
{\tiny\seqsplit{[Cl][c]1CCCOS[c](c2=N[C@@H]([CH2]CC(=O)[NH][CH2]C)c3[n][n]c([CH3])C3-c3nc2[cH]c(S[CH3])[cH]C3)[cH][cH]1}}
&
7.9336 & 0.8205 & -8.0426 & 0.2273 & 6.5092 \\
\hline
%       -- Molecule 4       --
\begin{adjustbox}{max width=2.4cm, max height=2.4cm, center}
\includegraphics{composite_cnn/4.png}
\end{adjustbox}
&
{\tiny\seqsplit{[Cl]c1[cH][cH][c]([C]2=[N]c([CH2][C](=[O])[NH][CH2][CH3])[c]3[n][n][c]([CH3])c3-[c]3[c]2[cH][c](C[CH3])[cH]c3)[cH][cH]1}}
&
7.0626 & 0.9165 & -8.4464 & 0.3023 & 6.4728 \\
\hline
%       -- Molecule 5       --
\begin{adjustbox}{max width=2.4cm, max height=2.4cm, center}
\includegraphics{composite_cnn/5.png}
\end{adjustbox}
&
{\tiny\seqsplit{[n]1[c]([CH3])cc2[c](n1)[C@H]([CH2][C](Cl)=O)[N]=[C]([c]1[cH][cH][c](CO)n[cH]1)[c]1[cH][cH][cH]c[c]1-2}}
&
6.9081 & 0.9228 & -8.8071 & 0.3378 & 6.3748 \\
\hline
\end{tabular}
\end{adjustbox}
\end{table*}

 
% ================================================================
%  FIGURE CAPTION — Top-5 Strategy 1 molecules
%  Paste as \caption{} or as legend paragraph after \end{figure}
% ================================================================
 
\noindent\small
\textbf{Table~1~$|$~Top-5 BRD4 candidate molecules identified by
non-interaction-site masking and Tanimoto-cutoff sequential
filtering (Strategy~1).}
Molecules were generated by progressively masking atoms outside
the PLIP-identified pharmacophoric interaction sites, allowing
ChemBERTa to reconstruct structurally varied analogues while
preserving the pharmacophore as unmasked context.
The top-5 candidates were selected from the docked library by
applying a four-step sequential filter (Algorithm~\ref{alg:idea1}):
CNNscore~$> 0.80$, Vinardo score~$< -8.0$~kcal/mol, Tanimoto
similarity to the reference ligand~$< 0.40$, and ranking by
composite score $S = \text{CNNaffinity} \times \text{CNNscore}$
(descending).
All five molecules satisfy the scaffold-novelty criterion
(Tanimoto~$< 0.40$) and carry a para-chlorophenyl moiety
together with a triazole or pyrazole nitrogen-containing core
consistent with the acetyl-lysine mimicry pharmacophore of
BRD4 bromodomain inhibitors~\cite{filippakopoulos2010selective}.
Composite scores range from $S = 6.77$ (Rank~1;
CNNaffinity~$= 7.40$, CNNscore~$= 0.91$,
Vinardo~$= -9.06$~kcal/mol) to $S = 6.37$ (Rank~5),
and all five Vinardo scores fall below $-8.0$~kcal/mol,
indicating strong predicted physical binding to the BRD4
pocket consistent with known high-affinity BET
inhibitors~\cite{mcnunn2021gnina}.
The full top-10 ranked list is provided in
Supplementary Table~S1.
\normalsize

 \begin{table*}[ht]
\centering
\caption{Butina scaffold-cluster representative selection}
\label{tab:butina_representatives}
\renewcommand{\arraystretch}{1.3}
\begin{adjustbox}{width=\textwidth}
\begin{tabular}{|
>{\centering\arraybackslash}m{2.6cm}|
>{\raggedright\arraybackslash}p{6.2cm}|
c|c|c|c|c|}
\hline
\textbf{Structure} &
\textbf{SMILES} &
\shortstack{\textbf{CNNaffinity}\\(higher is better)} &
\shortstack{\textbf{CNNscore}\\(higher is better)} &
\shortstack{\textbf{Vinardo}\\(lower is better)} &
\shortstack{\textbf{Tanimoto}\\(lower = more novel)} &
\shortstack{\textbf{Composite}\\(aff $\times$ score)} \\
\hline

%       -- Molecule 1       --
\begin{adjustbox}{max width=2.4cm, max height=2.4cm, center}
\includegraphics{vinardo_docking/top_molecules/1.png}
\end{adjustbox}
&
{\tiny\seqsplit{[CH3]Cc1[n][n][c]2[n]1-[c]1[c]([c]([CH3])[c]([CH3])[s]1)[C]([c]1[cH][cH][c]([Cl])[cH][cH]1)=[N][C@H]2[CH2]C(=[O])C[C]([CH3])([CH3])[CH3]}}
&
7.5949 & 0.9707 & -9.1590 & 0.6338 & 7.3724 \\
\hline
%       -- Molecule 2       --
\begin{adjustbox}{max width=2.4cm, max height=2.4cm, center}
\includegraphics{vinardo_docking/top_molecules/2.png}
\end{adjustbox}
&
{\tiny\seqsplit{[Cl][c]1CCcS[c](C2=N[C@@H]([CH2]C(=O)[NH][CH2]C)c3[n][n]c([CH3])N3-c3c2[cH]c(C[CH3])[cH]c3)[cH][cH]1}}
&
7.5955 & 0.9514 & -9.4728 & 0.5301 & 7.2265 \\
\hline
%       -- Molecule 3       --
\begin{adjustbox}{max width=2.4cm, max height=2.4cm, center}
\includegraphics{vinardo_docking/top_molecules/3.png}
\end{adjustbox}
&
{\tiny\seqsplit{[CH3]c1[n][n][c]2[n]1-[c]1[c](c(Cl)[c](Cl)[s]1)[C]([c]1[cH][cH][c]([Cl])[cH][cH]1)=[N][C@H]2N[C](=[O])[O][C]([CH3])([CH3])[CH3]}}
&
7.4721 & 0.9641 & -9.0129 & 0.5616 & 7.2041 \\
\hline
%       -- Molecule 4       --
\begin{adjustbox}{max width=2.4cm, max height=2.4cm, center}
\includegraphics{vinardo_docking/top_molecules/4.png}
\end{adjustbox}
&
{\tiny\seqsplit{[CH3]CCCc1[n][n][c]2[n]1-[c]1[c]([c]([CH3])[c]([CH3])[s]1)[C]([c]1[cH][cH][c]([Cl])[cH][cH]1)=[N][C@H]2[CH2]S(=[O])C[C]([CH3])([CH3])[CH3]}}
&
7.6885 & 0.9293 & -8.8238 & 0.5325 & 7.1450 \\
\hline
%       -- Molecule 5       --
\begin{adjustbox}{max width=2.4cm, max height=2.4cm, center}
\includegraphics{vinardo_docking/top_molecules/5.png}
\end{adjustbox}
&
{\tiny\seqsplit{[Cl]CCN1[cH][cH][c]([C]2=[N][C@@H]([CH2][C](=[O])[NH][CH2][CH3])[c]3[n][n][c]([CH3])[n]3-[c]3[c]2[cH][c]([O][CH3])[cH][cH]3)cc[cH]1}}
&
7.3375 & 0.9704 & -8.4429 & 0.6538 & 7.1201 \\
\hline
\end{tabular}
\end{adjustbox}
\end{table*}
 
% ===============

% ================================================================
%  FIGURE CAPTION — Top-5 Strategy 5 molecules
%  Paste as \caption{} or as legend paragraph after \end{figure}
% 
\FloatBarrier
\noindent\small
\textbf{Table~2~$|$~Top-5 scaffold-diverse BRD4 candidate
molecules identified by non-interaction-site masking and
Butina scaffold-cluster representative selection (Strategy~2).}
Molecules were generated by masking non-interaction atoms and
reconstructed by ChemBERTa, as described above.
Scaffold-level diversity was enforced by applying the Butina
single-linkage clustering algorithm~\cite{butina1999}
(Algorithm~\ref{alg:idea5}) with a Tanimoto distance cutoff
$\varepsilon = 0.40$ to partition the full docked library
into chemically distinct scaffold clusters, then selecting
the highest-composite representative from each cluster.
Unlike Strategy~1, no fixed Tanimoto cutoff to the reference
ligand is imposed; instead, diversity is guaranteed across
the selected set by construction, as no two candidates
belong to the same scaffold cluster.
The top-5 representatives span distinct chemical scaffolds
and collectively cover a broader region of BRD4-relevant
chemical space than any single-criterion selection strategy,
making them particularly suited to downstream
structure--activity relationship (SAR) analysis across
multiple scaffold classes~\cite{bemis1996scaffolds}.
All five candidates achieve CNNscore~$> 0.80$ and
Vinardo score~$< -8.0$~kcal/mol, confirming both
pose quality and physical binding plausibility
within the BRD4 acetyl-lysine pocket.
The full top-10 scaffold-diverse candidate list is provided
in Supplementary Table~S2.
\normalsize

\clearpage 
\section{Discussion}


 
The comparative analysis across Figures 2–4 demonstrates that non-interacting part (NIP Masking ) masking systematically altered the generative behavior of ChemBERTa in a scaffold-dependent manner, affecting not only the number of valid molecules generated but also their docking quality, novelty, and internal chemical diversity.

For 1GH-A-1173, NIP Masking  produced the clearest advantage over random masking. Figure 2 shows that the majority of valid molecules were recovered at very low masking depths, with the generation curve peaking sharply at mask count = 2, whereas random masking collapsed rapidly after mask count = 3. This indicates that preserving PLIP-identified pharmacophoric interaction centers while perturbing peripheral non-binding regions provides sufficient structural flexibility for ChemBERTa to explore nearby chemical space without disrupting the core BRD4 recognition motif. Importantly, this reduced masking requirement implies lower computational cost because productive exploration occurred before extensive token masking became necessary. The effect is particularly significant for compact scaffolds such as 1GH, where masking interaction-critical atoms likely destroys a disproportionate amount of chemically informative context.

The superiority of NIP Masking  for 1GH extended beyond generation yield into novelty and docking quality. In Figures 3 and 4, NIP Masking  generated approximately five times more dockable molecules than random masking and substantially increased the number of novel-scaffold candidates in the Tanimoto$<$ 0.40 region. Moreover, these novel molecules retained competitive composite docking scores and Vinardo affinities, indicating that structural novelty was not achieved at the expense of predicted binding quality. The pairwise similarity analysis in Figure 5 further supports this observation: the NIP Masking library displayed markedly lower mean pairwise similarity than the random library, demonstrating substantially greater internal diversity. Collectively, these findings suggest that preserving interaction hotspots stabilizes the pharmacophoric architecture while allowing peripheral exploration of chemically diverse substituents.

EAM-A-1 exhibited a similarly strong NIP Masking advantage, although with a different diversity profile. NIP Masking  consistently generated more molecules across all masking depths and maintained productive generation at larger mask counts where random masking rapidly decayed. In the docking analysis, EAM produced the highest composite docking scores observed in the study, and high-scoring molecules emerged early in the generation process. Biochemically, this behavior is likely related to the rigid tricyclic benzodiazepine core and the chlorine-mediated hydrophobic contact identified by PLIP. The para-chlorophenyl region probably functions as a strong hydrophobic anchor within the BRD4 acetyl-lysine pocket, meaning that preserving these interaction-rich regions maintains favorable enthalpic contacts during generation. However, unlike 1GH, EAM showed relatively lower similarity to the original ligand under NIP Masking, indicating that the model successfully shifted away from the starting structure to explore novel chemical space. This suggests that while 1GH's rigid privileged scaffold constrains ChemBERTa toward local, potency-preserving analogues, EAM allows for much broader scaffold diversification.

In contrast, 62G-A-201 represented the principal exception to the overall trend. Random masking generated substantially more dockable molecules than NIP Masking  in both the generation curves and docking analyses. Nevertheless, the qualitative distribution of docking scores revealed an important distinction: random masking produced many low-quality binders, whereas NIP Masking  concentrated molecules within the stronger CNNaffinity and Vinardo ranges. This indicates a quantity-versus-quality trade-off. A likely explanation is the chemically tolerant benzodiazepine–isoxazole topology of 62G. The scaffold contains multiple aromatic and hydrophobic regions that can accommodate substitution without catastrophic loss of BRD4 binding. As a result, random perturbations remain chemically viable more often than in compact interaction-dense systems such as 1GH. Additionally, benzodiazepine-like chemotypes are heavily represented in PubChem and ChEMBL, the same large-scale corpora used during ChemBERTa pretraining. The pretrained model therefore likely learned robust token-level priors for these motifs, allowing successful reconstruction even under random masking.

The comparison between JQ1-A-201 and JQ1-A-1 further demonstrates how subtle interaction differences reshape the generative landscape. Although both structures share the same thienodiazepine-triazole scaffold, their PLIP interaction maps differ: JQ1-A-201 contains water-bridge-mediated interactions, whereas JQ1-A-1 replaces these with a direct hydrogen bond. These differences alter which molecular regions remain preserved during masking and consequently change the chemical neighborhoods explored by ChemBERTa. For JQ1-A-1, random masking generated more dockable molecules and discovered a larger number of unique molecules at low masking depths, especially at mask counts 2 and 3. This suggests that the JQ1 scaffold possesses several chemically permissive peripheral positions where perturbation does not severely disrupt BRD4 recognition. Conversely, JQ1-A-201 exhibited more balanced behavior between strategies, with NIP Masking  preferentially enriching molecules in the scaffold-transition region near the Tanimoto novelty boundary.

Across all ligand systems, the combined analysis of CNNscore, CNNaffinity, and Vinardo scoring demonstrates that molecule quantity alone is not an adequate measure of generative success. Random masking frequently increased the total number of generated molecules but also increased the proportion of weak binders and high-similarity analogues. In contrast, NIP Masking  generally produced more coherent chemical neighborhoods, stronger docking-score distributions, and improved novelty–potency balance. The Tanimoto analyses further support this interpretation by showing that NIP Masking  more consistently generated unimodal similarity distributions, whereas random masking often produced bimodal or heavy-tailed profiles enriched with near-identical analogues. This suggests that preserving interaction-critical regions biases ChemBERTa toward chemically meaningful exploration trajectories rather than unconstrained stochastic perturbation.

Overall, the results demonstrate that preserving non-covalent interaction hotspots during masking can guide ChemBERTa toward chemically coherent and pharmacologically plausible exploration trajectories. However, the effectiveness of NIP Masking  is highly scaffold-dependent and reflects an interplay between binding-pocket interaction topology, scaffold rigidity, pharmacophoric redundancy, and biases inherited from large-scale chemical pretraining datasets.

To contextualise the generated molecules within the broader landscape
of known bioactive compounds, each library was compared against the
ChEMBL database using nearest-neighbour Tanimoto similarity
(Figs.~\ref{fig:nn_similarity_hist}--\ref{fig:boxplot_chembl_summary}).
The results consistently show that NIP masking and random masking differ
not only in how many molecules they generate, but in where those molecules
sit relative to known pharmacological chemical space.
In four of the five systems, NIP masking produced nearest-neighbour
distributions centred in the mid-similarity range ($T = 0.40$--$0.70$),
whereas random masking shifted distributions toward the high-similarity
tail ($T \geq 0.80$), indicating a tendency to recover near-identical
ChEMBL compounds rather than structurally distinct analogues.
The difference was most pronounced for JQ1-A-1
(Fig.~\ref{fig:nn_similarity_hist}g,h), where random masking
produced a strongly right-skewed distribution with a sharp peak
at $T \approx 0.80$, while NIP masking generated a unimodal
distribution centred at $T \approx 0.60$ in the medium-similarity
range, and for 1GH-A-1173 (Fig.~\ref{fig:nn_similarity_hist}a,b),
where NIP masking simultaneously covered a novel-scaffold peak near
$T \approx 0.30$ and a broader analogue population at
$T = 0.45$--$0.65$, while the random library ($n = 21$) was
concentrated at $T \geq 0.75$ and largely absent from both of
these chemically relevant regions.
The frequency analysis in Figure~\ref{fig:nn_frequency} further shows
that NIP masking libraries mapped to a larger number of distinct ChEMBL
compounds across most systems: for JQ1-A-1, NIP masking identified
nine secondary ChEMBL neighbours beyond the dominant hit
(CHEMBL1957266, $n \approx 230$, Med:~$0.65$) compared to only five
under random masking ($n \approx 135$, Med:~$0.70$), and for EAM-A-1,
NIP masking mapped molecules to eight distinct ChEMBL compounds versus
five for random, with a lower median similarity to the top hit
(CHEMBL1232461, PDB:~7AQT; Med:~$0.74$ vs.\ $0.85$), indicating
greater structural variation within the NIP masking library.
Importantly, the dominant ChEMBL nearest neighbours were
pharmacologically interpretable across all five systems:
CHEMBL2181820 for 1GH-A-1173 (Med:~$0.62$),
CHEMBL1957266 (the JQ1 parent compound) for both JQ1 crystal forms,
CHEMBL4303404 (PDB:~5HLS, the reference compound itself) for
62G-A-201, and CHEMBL1232461 (PDB:~7AQT) for EAM-A-1, confirming
that both strategies correctly anchor generation within BRD4-relevant
chemical space.
EAM-A-1 was again the exception: NIP masking produced a tightly
clustered distribution at high ChEMBL similarity
($T \approx 0.75$--$0.80$; Fig.~\ref{fig:nn_similarity_hist}i),
consistent with the rigid tricyclic benzodiazepine core constraining
reconstruction to a narrow chemical neighbourhood, whereas random
masking generated a bimodal distribution spanning a novel-scaffold
region near $T \approx 0.30$ and a high-similarity peak near
$T \approx 0.85$, suggesting that random perturbations of the rigid
scaffold sample two disconnected regions of ChEMBL chemical space.
The cross-ligand boxplot summary (Fig.~\ref{fig:boxplot_chembl_summary})
confirms this pattern at the aggregate level: NIP masking yields
tighter, more consistent nearest-neighbour distributions across all
five systems, with all median values above the $T = 0.40$
pharmacological relevance threshold, demonstrating that
interaction-guided generation produces libraries that are both
chemically novel and pharmacologically grounded.

\subsubsection{Non Interacting part masking makes molecules close to ChEMBL active BRD4 inhibitors}


The nearest-neighbour frequency analysis in Figure~7 identifies which specific ChEMBL compounds (7852 Known active BRD4 binders: In the GenPLIP analysis, "known active BRD4 binder" refers to a small molecule compound from ChEMBL that has been experimentally confirmed to bind to the BRD4 protein with measurable affinity above a defined potency threshold) the generated libraries cluster around and how structurally diverse those clusters are across the two masking strategies.  Across four of the five systems, NIP masking produced a larger panel of
distinct ChEMBL nearest neighbours than random masking, indicating that interaction-guided generation steers ChemBERTa toward multiple
pharmacologically validated scaffolds rather than concentrating the
library around a single dominant hit.







For the 1GH-A-1173 ligand system(Fig.~7a,b), NIP masking generated molecules 
clustering around two pharmacologically validated BRD4-relevant 
\texttt{ChEMBL} neighbours  predominantly \texttt{CHEMBL2017291} 
(I-BET151/GSK1210151A), a well-characterised BET bromodomain inhibitor 
with 688 bioactivity records dominated by IC\textsubscript{50} and 
general activity measurements, and secondarily \texttt{CHEMBL2017287}, 
a minimally characterised epigenetic regulator ligand with only 9 
bioactivity records where general Activity and binding affinity (Kd) 
measurements predominate with both neighbours falling in the 
medium-to-high similarity tier, indicating pharmacologically grounded 
yet a structurally diverse generation. Random masking, by contrast, spread 
its output across three neighbours with \texttt{CHEMBL2017287} becoming 
the dominant hit, demoting the well-validated \texttt{CHEMBL2017291} to 
a secondary role, and additionally recovering \texttt{CHEMBL4862592}, a 
minimally profiled compound with only 2 IC\textsubscript{50} records 
against a single epigenetic regulator target. Taken together, these results confirm that NIP masking produces a more pharmacologically 
anchored library, correctly prioritising the well-validated BET 
bromodomain inhibitor as the dominant nearest neighbour, whereas random 
masking disperses generation toward more poorly characterised chemicals 
space. The two-dimensional structures of all \texttt{ChEMBL} nearest neighbours identified in this analysis are provided in the Supplement~3.


For the 62G-A-201 ligand system (Fig.~7e,f), NIP masking generated a substantially 
richer chemical neighbourhood than random masking, mapping molecules to 
eight distinct \texttt{ChEMBL} neighbours compared to only five under 
random masking. In both strategies, \texttt{CHEMBL4303404} (5HLS) a 
moderately characterised BRD4-relevant compound with 131 bioactivity 
records dominated by IC\textsubscript{50} and activity measurements 
across epigenetic regulators and enzyme targets  emerged as the 
overwhelmingly dominant nearest neighbour, confirming that both 
strategies correctly anchor generation to this pharmacologically 
validated reference scaffold. Both strategies also recovered 
\texttt{CHEMBL5274793}, \texttt{CHEMBL5280537}, and \texttt{CHEMBL5266575} 
  all sparsely characterised ADME-focused compounds with 10 bioactivity 
records dominated by CL and PPB measurements against the epigenetic regulator 
targets. NIP masking exclusively recovered four additional neighbours 
absent from random masking: \texttt{CHEMBL2431081}, \texttt{CHEMBL5276062}, 
and \texttt{CHEMBL5905822} all minimally profiled compounds with 2--5 
bioactivity records characterised via IC\textsubscript{50} and kinetic 
binding parameters against single epigenetic regulator targets and 
\texttt{CHEMBL5289896}, the only high-similarity neighbour in this system. 
Random masking exclusively recovered \texttt{CHEMBL5278351}, a moderately 
characterised ADME compound with 20 bioactivity records. Taken together, 
These results quantitatively confirm that masking the non-interacting 
components successfully align zero-shot generative outputs with 
pharmacologically validated chemical space, yielding a broader panel of 
distinct \texttt{ChEMBL} nearest neighbours rather than over-concentrating the generated library around a single reference scaffold. 


For JQ1-A-1 (Fig.~7g,h),  NIP masking identified seven secondary ChEMBL
neighbours versus five for random masking beyond the dominant hit in both masking strategies.
\texttt{CHEMBL1957266} (active enantiomer of JQ1). \texttt{CHEMBL1957266} [(+)-JQ1], the benchmark BET inhibitor used as the reference active compound, exhibits a uniquely broad assay footprint in ChEMBL spanning direct binding, functional cellular readouts, and pan-BET epigenetic regulator engagement. Unlike the broader BRD4 dataset, where ~95 percent of assays were direct binding (B-Binding), here F-Functional assays constitute a substantial fraction (~35 percent), reflecting JQ1's widespread use in cell-based cancer proliferation, transcription, and gene-expression studies beyond simple biochemical binding confirmation. In NIP masking, more than 130 molecules have  \texttt{CHEMBL1957266} as nearest neighbour, with 50 molecules having higher similarity, and random masking has 165 molecules but with low similarity. \texttt{CHEMBL4749839} has structure is identical to JQ1 (same thienotriazolodiazepine scaffold, 4-chlorophenyl, three methyl groups) except the tert-butyl ester (–COOC(CH₃)₃) of JQ1 is replaced by a tert-butyl ketone (–COC(CH₃)₃), reducing the formula from C₂₃H₂₅ClN₄O₂S (JQ1, MW 457) to C₂₃H₂₅ClN₄OS (MW 441). This single-oxygen removal eliminates the ester's metabolic liability — JQ1's notoriously short in vivo half-life (~1 h) is partly attributed to ester hydrolysis — making \texttt{CHEMBL4749839} a metabolically stabilised, preclinical-stage JQ1 derivative only having rediscovered by NIP masking, indicating random masking failed to rediscover important molecules and keeping the interaction part fixed was helpful for ChemBerta inference. \texttt{CHEMBL3769755} (PDB id: 6DUV) is present in both masking  strategies, the most polar and soluble member of the JQ1 ester series, with a decade of sustained research interest
The structure is JQ1 with the tert-butyl ester replaced by a methyl ester (–CH₂COOCH₃), shrinking the formula from C₂₃H₂₅ClN₄O₂S (JQ1, MW 457) to C₂₀H₁₉ClN₄O₂S (MW 414, a loss of C₃H₆ = the three tert-butyl carbons), making it the smallest, most hydrophilic ester variant of the JQ1 scaffold; its 10 publications spanning continuously from 2015 to 2024 with renewed activity in 2024 — confirms it as a sustained reference compound in BRD4 medicinal chemistry, not a transient screening hit Both strategy found this with NIP molecules having slightly higher similarity meaning they it is heavily present in chemberta. Molecules like \texttt{CHEBML518340}  (a direct ester (nor-methylene) analogue of JQ1 — the CH₂ spacer is surgically removed). \texttt{ChEMBL2349340} is a preclinical-stage compound with a molecular weight of 456.01, featuring a thienotriazolodiazepine-like core structure with a tert-butyl ester group. \texttt{CHEBML518340} and  \texttt{CHEMBL2349340} were found in both strategies. \texttt{CHEMBL3967931}, an epigenetic regulator target only found in NIPS masking. Showing that NIP masking's superiority and diversity in ChemBL dataset exporation then random masking.  

NIP masking has a higher similarity with
\texttt{CHEMBL5181340}, \texttt{CHEMBL3769755}
(6DUV) Then random masking.     Moreover,  \texttt{CHEMBL4749839},\texttt{CHEMBL5822147} and \texttt{CHEMBL3967931} were found only in NIP Masking similarity matching.
JQ1-A-1, NIP masking not only found more molecules in similarity search eight versus random four, but random masking also completely failed to get molecules generated similar to PDB id 6DUV (\texttt{CHEMBL3769755}).

%             -
For JQ1-A-201 (Fig.~7c,d), the same trend observed for JQ1-A-1 holds, NIP masking recovers a markedly richer set of ChEMBL nearest neighbours than random masking of eight reference compounds (seven secondary neighbours beyond the dominant hit) versus only four (three secondary neighbours).
Under both strategies, the dominant neighbour is \texttt{CHEMBL1957266} [(+)-JQ1], the benchmark BET inhibitor described above. The two strategies differ in match quality rather than in identity. NIP masking maps ~133 generated molecules onto (+)-JQ1 with a median Tanimoto of 0.62, so most of the population sits in the medium to high similarity range. Random masking maps more molecules (~148) but at a lower median of 0.58, with the majority falling in the low band (0.40–0.60). The high similarity fractions are comparable (~38 molecules each). NIP, therefore, concentrates generation around (+)-JQ1 at genuinely higher chemical similarity, whereas random masking inflates the raw count with weakly related structures, the same, fewer but closer (NIP) versus more-but-weaker (random) behaviour seen for JQ1-A-1.

The advantage of NIP is clearest among the secondary neighbours. The two chemically most informative analogues are recovered only by NIP \texttt{CHEMBL4749839} (median 0.60, ~11 molecules), the tert-butyl-ketone analogue of JQ1 described above, and \texttt{CHEMBL3769755} (6DUV median 0.52, ~4 molecules), the methyl ester variant described above. Random masking fails to rediscover either most notably 6DUV, the sustained reference compound in BRD4 medicinal chemistry, reproducing the JQ1-A-1 finding that keeping the interaction region fixed helps ChemBERTa retrieve the relevant JQ1 chemical series. Three further neighbours appear exclusively under NIP \texttt{CHEMBL5851164} (median 0.56), \texttt{CHEMBL5822147} (median 0.44) and \texttt{CHEMBL3967931} (median 0.71, the single highest-similarity secondary neighbour of the panel), the latter two also being NIP exclusive for JQ1-A-1.
Two neighbours are shared by both strategies, and here the comparison is more nuanced. For \texttt{CHEMBL2349340}, the two are close, with NIP marginally higher in similarity (0.54 vs 0.49) and random marginally higher in count (~4 vs ~2). For \texttt{CHEMBL5181340}, however, the ordering reverses relative to JQ1-A-1 random masking matches it at higher similarity (0.68, medium band) than NIP (0.46, low band) at comparable counts (~12 vs ~14). Random masking also contributes one neighbour of its own, \texttt{CHEMBL3889858} (median 0.51, a single molecule), absent from the NIP set.

Taken together, JQ1-A-201 reinforces the JQ1-A-1 conclusion NIP masking both broadens the neighbour set (eight versus four) and recovers the reference active and its key preclinical/structural analogues at higher similarity, while random masking completely fails to regenerate the JQ1 ester series exemplified by 6DUV. The two exceptions, random's higher similarity match to \texttt{CHEMBL5181340} and its unique recovery of \texttt{CHEMBL3889858}, show that the non-interaction fixed (NIP) representation is not strictly dominant for every individual neighbour, but it is decisively better at recovering the chemically meaningful JQ1 scaffold variants.

%             - 
For EAM-A-1 (Fig.~7i,j), NIP masking identified four secondary ChEMBL neighbours
versus three for random masking, beyond the dominant hit common to both masking
strategies. \texttt{CHEMBL1232461} (7AQT) molibresib (I-BET762\,/\, GSK525762), the
benchmark BET inhibitor used as the reference active compound, and itself the BET
ligand bound to the EAM-A-1 target, BRD4-BD2 (PDB 7AQT), exhibits a uniquely
broad assay footprint in ChEMBL, spanning direct binding (B-Binding), functional
cellular readouts (F-Functional), and pan-BET epigenetic regulator engagement.
This reflects its dual role as a clinical BET candidate (e.g.,~NUT-midline
carcinoma) and a widely used chemical probe, so that F-Functional assays form a substantial fraction of its record rather than the binding-dominated profile
typical of the broader BRD4 dataset. In NIP masking, $\sim$107 generated molecules
have \texttt{CHEMBL1232461} as nearest neighbour (median Tanimoto~$\approx 0.68$), the bulk
falling in the medium similarity bin (0.60--0.80) with a sizeable low similarity
tail (0.40--0.60, $\sim$30 molecules); random masking returns a near identical
$\sim$105 molecules (median~$\approx 0.69$) but shifted slightly toward the medium
bin with a much smaller low similarity tail ($\sim$10 molecules). The two shared
secondary neighbours, \texttt{CHEMBL5898140}(median~$\approx 0.75$ in both) and
\texttt{CHEMBL5931751} ($\approx 0.64$ NIP\,/\,$\approx 0.65$ random), behave
near identically across strategies. Beyond these, \texttt{CHEMBL6040604} and \texttt{CHEMBL5917253}
(medians~$\approx 0.45$ and~$\approx 0.41$) were recovered only under NIP masking,
whereas \texttt{CHEMBL1958347} (median~$\approx 0.43$) appeared only under random masking.
Thus, for EAM-A-1, NIP masking recovered a marginally broader secondary neighbour
set (four versus three) and uniquely captured \texttt{CHEMBL6040604} and \texttt{CHEMBL5917253},
which random masking failed to generate; the dominant molibresib cluster, however,
was essentially unchanged in size ($\sim$107 vs.\ $\sim$105) and median similarity
($\approx 0.68$ vs.\ $\approx 0.69$), and no secondary neighbour reached systematic
high similarity ($T \geq 0.80$) under either masking scheme.
%             -

Importantly, the dominant ChEMBL nearest neighbour was pharmacologically
interpretable in every case: \texttt{CHEMBL2181820} for 1GH-A-1173 (Med: 0.62),
the JQ1 parent compound \texttt{CHEMBL1957266} for both JQ1 crystal forms, the
reference compound \texttt{CHEMBL4303404} (PDB: 5HLS) for 62G-A-201, and
\texttt{CHEMBL1232461} (PDB: 7AQT aka MOLIBRESIB a Phrase 2 drug) for EAM-A-1.
The consistent recovery of these known BRD4-active scaffolds as dominant
nearest neighbours confirms that both strategies remain pharmacologically
anchored to the correct target class; what distinguishes NIP masking is
that it reaches those anchors from a broader and more structurally
diverse set of generated molecules.
 
The colour composition of the bars reinforces this interpretation.
Under NIP masking, the dominant bar for  is predominantly
darker green then random masking  indicating that the majority of
generated molecules are related to known BRD4 inhibitors.


In brief, preserving PLIP-identified interaction hotspots
during masking biases ChemBERTa toward the pharmacophore in ways that
allow peripheral structural variation, whereas random masking disrupts
chemically informative context and forces the model to rely on global
sequence statistics, which are inherently biased toward high-frequency,
high-similarity motifs in the pre-training corpus.
 

 
Whereas Figure~7 characterises which specific ChEMBL compounds are
matched and by how many molecules, the aggregate boxplot in Figure~8
answers a different but complementary question: how \emph{consistently}
does each masking strategy maintain pharmacological relevance across
structurally diverse scaffolds?
Each box in Figure~8 summarises the full distribution of
per-molecule nearest-neighbour Tanimoto scores for one ligand group;
the red line marks the group median; the box edges span the
interquartile range (IQR); and circles denote outliers beyond
$1.5 \times \text{IQR}$.
The grey dashed reference at $T = 0.40$ marks the conventional
scaffold-similarity threshold~\cite{maggiora2014} below which a generated molecule
has no recognised bioactive analogue in ChEMBL~\cite{gaulton2012chembl}.
The NIP masking generated molecules are consitantly more similar to ChEMBL for all five ligands. 

The combined picture presented by Figures~7 and~8 carries a consequence
that extends beyond the immediate comparison of masking strategies.
A generative model that reliably places its output in the
$T \in [0.40, 0.80)$ analogue zone of known bioactive compounds does
not merely score well on a benchmarking metric; it produces molecules
with a meaningfully higher probability of success at every subsequent
stage of the drug discovery pipeline.
ChEMBL entries in this similarity range have been shown to be
synthetically accessible, to possess physicochemical profiles compatible
with oral bioavailability, and to carry at least partial experimental
validation of target engagement~\cite{gaulton2012chembl}.
A generated molecule anchored to such a neighbour therefore inherits
that body of precedent: synthetic routes are likely to be adaptable
from known chemistry, ADMET liabilities are broadly predictable from
the known analogue's profile, and the pharmacophoric features
responsible for BRD4 recognition have already been validated in
biochemical assays.
This is the lead-hopping region of chemical space novel enough to
offer new intellectual property and potentially improved potency or
selectivity, yet close enough to established actives to be optimised
with confidence rather than requiring first-principles characterisation
of an entirely new scaffold.
 
Equally important is what Figures~7 and~8 reveal about structural
diversity \emph{within} that pharmacologically grounded space.
A library that clusters exclusively around a single dominant ChEMBL hit
at high median similarity ($T \geq 0.80$) offers little scope for
structure--activity relationship (SAR) exploration: all candidates share
essentially the same scaffold, so any common liability metabolic
instability, off-target activity, poor solubility will affect the
entire set simultaneously~\cite{Paul2010}.
The broader ChEMBL coverage demonstrated by NIP masking in Figure~7,
where generated molecules map to multiple distinct known BRD4 inhibitors
at lower median similarity, means that the NIP library spans several
pharmacophore realisations simultaneously.
This is precisely the portfolio composition required for productive
hit-to-lead campaigns: multiple structurally distinct series can be
advanced in parallel, providing redundancy against scaffold-specific
attrition and enabling cross-series SAR comparisons that identify which
structural features are essential for potency and which are tolerant of
modification~\cite{bemis1996scaffolds}.
The tight, consistent IQR distributions shown for NIP masking in
Figure~8 add a further practical dimension: a medicinal chemist
synthesising from this library can expect predictable structural
relationships between members, making iterative analoguing more
tractable than if the library exhibited the wide molecule-to-molecule
variability seen under random masking~\cite{maggiora2014}.
Taken together, the results of Figures~7 and~8 therefore demonstrate
that GenPLIP's interaction-aware masking strategy does not merely
generate chemically valid SMILES or improve docking scores in isolation;
it produces libraries whose position in chemical space is simultaneously
pharmacologically precedented, structurally diverse across multiple
BRD4-relevant chemotypes, and internally consistent enough to support
the kind of systematic SAR analysis that translates computational hits
into viable drug candidates.

\subsubsection{Top candidates based on novelty potency analysis}
Tables~1 and~2 gives us based on novelty potency analysis but also answer in two distinct ordering from the same
NIP masking-generated library: Table~1 asks which candidates
are most structurally novel relative to the reference ligand,
while Table~2 asks which candidates collectively cover the
broadest range of distinct scaffolds within the library itself.
Both are necessary because novelty against a single reference
compound and diversity within a selected set are independent
properties. A molecule that is highly novel relative to the
reference ligand may be nearly identical to another candidate
in the same selected set, and a set optimised purely for
within-set diversity may contain members that closely resemble
the starting compound.
Together, Tables~1 and~2 therefore present the top-5 BRD4
candidate molecules from the NIP masking library under two
complementary selection strategies, demonstrating that
interaction-aware generation can yield structurally novel,
high-affinity candidates that satisfy all standard
medicinal-chemistry quality criteria simultaneously.
 
The five candidates in Table~1, selected by Tanimoto-cutoff
sequential filtering (Strategy~1), collectively establish that
NIP masking produces genuine scaffold-novel BRD4 inhibitors
rather than minor modifications of the reference ligand.
All five satisfy the scaffold-novelty criterion
(Tanimoto $< 0.40$ to the reference ligand~\cite{maggiora2014}),
confirming that the generated structures occupy a distinct
region of chemical space from the co-crystal starting point,
yet all retain composite scores $S = \text{CNNaffinity} \times
\text{CNNscore}$ ranging from $S = 6.77$ (Rank~1;
CNNaffinity $= 7.40$, CNNscore $= 0.91$,
Vinardo $= -9.06$~kcal/mol) to $S = 6.37$ (Rank~5), and all
five Vinardo scores fall below $-8.0$~kcal/mol~\cite{quiroga2016vinardo},
confirming strong predicted physical binding to the BRD4
acetyl-lysine pocket.
Structurally, all five candidates carry a
\emph{para}-chlorophenyl moiety together with a triazole or
pyrazole nitrogen-containing core~\cite{filippakopoulos2010selective}, consistent with
the acetyl-lysine mimicry pharmacophore known to be essential
for BRD4 bromodomain recognition.
This convergence on a common pharmacophoric motif despite the
absence of any explicit pharmacophore constraint in the
generation pipeline is significant: it suggests that preserving
the PLIP-identified interaction hotspots during masking
implicitly encodes the BRD4 binding requirements into the
reconstruction context, guiding ChemBERTa toward
pharmacophorically coherent output without hard-coding
chemical rules.
The structural novelty of these candidates relative to the
reference ligand (Tanimoto range $0.23$--$0.40$) indicates
that the peripheral regions of each molecule have been
substantially redesigned, providing meaningful scope for
improved selectivity or pharmacokinetic properties over the parent
scaffold~\cite{mcnunn2021gnina}.
 
 
Whereas Table~1 optimises for absolute novelty relative to the
reference ligand, Table~2 addresses a different and
complementary objective: ensuring that the selected candidates
span the broadest possible range of distinct chemical scaffolds
present in the NIP masking library.
The five candidates selected by Butina scaffold-cluster
representative selection (Strategy~2) are guaranteed by
construction to belong to different scaffold clusters
($\varepsilon = 0.40$, consistent with the Bemis--Murcko
scaffold-equivalence threshold~\cite{bemis1996scaffolds}), meaning no two
candidates share the same core framework.
This scaffold-level diversity is reflected in the Tanimoto
values, which range from $0.53$ to $0.65$ higher than the
$< 0.40$ ceiling imposed in Strategy~1, indicating that some
Table~2 candidates are structurally closer to the reference
ligand, but the set as a whole covers a broader chemical
neighbourhood than any single-criterion approach could achieve.
Despite relaxing the novelty gate, all five Strategy~2
candidates achieve CNNscore $> 0.80$ and Vinardo score
$< -8.0$~kcal/mol, confirming that pose quality and physical
binding plausibility are maintained across all selected
scaffolds.
The top-ranked candidate achieves a composite score of
$S = 7.37$ (CNNaffinity $= 7.59$, CNNscore $= 0.97$,
Vinardo $= -9.16$~kcal/mol), which is notably higher than the
top Strategy~1 candidate ($S = 6.77$), reflecting the fact
that removing the hard Tanimoto novelty gate allows Strategy~2
to access higher-scoring molecules that would otherwise have
been filtered out.
This quality-versus-novelty trade-off is a deliberate design
feature: Strategy~2 is not intended to replace Strategy~1 but
to complement it by providing a set of structurally diverse
series suitable for parallel hit-to-lead development, where
each scaffold cluster can be independently optimised~\cite{bemis1996scaffolds}.
 
Comparing the actual results of the two tables side by side
makes these differences concrete.
In Table~1, all five candidates converge on a shared
pharmacophoric motif   a \emph{para}-chlorophenyl group
combined with a triazole or pyrazole nitrogen-containing
core   despite having Tanimoto similarities to the reference
ligand ranging from only $0.23$ to $0.40$, meaning the
peripheral regions of each molecule differ substantially while
the pharmacophore is preserved.
In Table~2, no such shared motif exists: the five candidates
belong to five distinct scaffold clusters and present
structurally different core frameworks, with Tanimoto
similarities to the reference ligand ranging from $0.53$ to
$0.65$.
For a medicinal chemist, this distinction is practically
significant: Table~1 provides one coherent scaffold series
that can be optimised by systematic peripheral substitution,
whereas Table~2 provides five separate starting points that
can be advanced in parallel, each representing a different
structural solution to BRD4 binding.
In terms of predicted binding quality, Table~2 candidates
achieve marginally stronger scores overall Vinardo scores
range from $-8.44$ to $-9.47$~kcal/mol compared with
$-8.04$ to $-9.20$~kcal/mol in Table~1, and CNNscores are
uniformly high in both sets ($> 0.82$ in Table~1,
$> 0.93$ in Table~2) reflecting the fact that removing
the hard novelty gate in Strategy~2 allows access to
higher-scoring molecules that Strategy~1 would exclude.
The key result is therefore not that one table is better than
the other, but that they capture different and complementary
regions of the quality--novelty--diversity space that the NIP
masking library spans.
 
\section{Conclusion}
The most immediate next step is experimental validation of the
highest-ranked GenPLIP candidates.
The top-scoring molecules from Tables~1 and 2 should be prioritised
for synthesis and \textit{in vitro} biochemical characterisation
against recombinant BRD4 bromodomain, followed by cellular
target-engagement assays and ADMET profiling.
Compounds that sustain potency through this preclinical cascade
represent genuine AI-generated BRD4 inhibitor leads, and the most
promising would be eligible for Investigational New Drug
(IND)-enabling studies and eventual Phase~I clinical evaluation a
milestone that would directly validate interaction-guided zero-shot
generation as a productive strategy at the hit-identification stage.
 
Looking further ahead, the NIP masking strategy is
architecture-agnostic and its application to other molecular
language models including SELFIES-based transformers, which
guarantee chemical validity by construction, and
autoregressive decoder models such as MolGPT and Chemformer,
which could treat interaction profiles as conditional tokens in
a sequence-to-sequence formulation represents a natural
extension for cross-model benchmarking.
Beyond adapting existing models, the most transformative
direction is a NIP-native expert language model pre-trained from
scratch on the full corpus of PLIP-annotated Protein Data Bank
co-crystal structures, with NIP masking as its sole training
objective.
Such a model would develop representations intrinsically organised
around non-covalent interaction geometry, generalise across target
classes, and serve as the foundation for closed-loop optimisation
in which experimentally measured binding affinities feed back as
reinforcement signals creating a self-improving generative system
that accumulates pharmacological expertise with each experimental
campaign.

















\newpage \section{References}
\vspace{-2em}
\bibliographystyle{unsrt}
\renewcommand{\refname}{}
\bibliography{references}

\end{document}

